% !TeX root = ../../Book2.tex
\chapter{变分法}
\label{b2:ch:14}
\BookChapterNumber{b2:ch:14}
% 本章摘要（不计入编号）
\begin{chapterabstract}
本章用函数空间中的微分研究积分泛函，推导标量与向量值 Euler--Lagrange 方程，并从边界项辨认自然边界条件。Dirichlet 能量说明凸性怎样把驻点变成唯一极小点；波动作用量则说明驻点未必是极小点。对梯度幂能量，直接方法通过强制性、弱闭约束与弱下半连续性构造极小元。若能量无下界，寻找驻点就需要另一条道路：选读部分用 Palais--Smale 条件控制近临界序列，以变形引理和山路定理寻找非极小型临界点，并应用于次临界半线性椭圆方程。
\end{chapterabstract}

% 本章目标（不计入编号）
\begin{learninggoals}
\begin{itemize}
\item\label{b2:item:14-goal-derivative}区分方向变分、Gâteaux 导数与 Fréchet 导数，检查积分下求导所需的控制。
\item\label{b2:item:14-goal-boundary}从允许扰动出发推导内部方程及自然边界条件，核对外法向的符号。
\item\label{b2:item:14-goal-energy}用精确能量差或第二变分判断极小性，区别作用量与守恒能量。
\item\label{b2:item:14-goal-direct}在 Sobolev 边界类中实施直接方法，指出强制性、弱紧性和下半连续性各自的作用。
\item\label{b2:item:14-goal-mountain}理解 Palais--Smale 条件控制什么序列，用下降变形解释山路临界点，并验证次临界椭圆模型中的几何条件与强收敛。
\end{itemize}
\end{learninggoals}

一条连接固定端点的曲线怎样使长度最小？一张固定在边框上的薄膜怎样达到平衡？这类问题比较的是整个函数，而不只是有限个参数。微积分仍然有效，但导数、紧性和边界条件都必须放到合适的函数空间中。候选曲线满足什么方程，最小值又是否真的能够达到，是两个相连的问题：前者由变分导数回答，后者需要研究极小化序列的极限。

\ProseParagraphBreak

极小化序列可能没有逐点收敛的子列，却仍能在弱收敛下保留能量信息。第一卷第九章的凸泛函理论解释这种极限为何有用；第二十一至二十四章的 Sobolev 空间与 Poincaré 不等式，则帮助我们同时控制函数、梯度和边界值。关于变分法的系统叙述，可进一步阅读 Bernard Dacorogna 的《Direct Methods in the Calculus of Variations》\cite{Dacorogna2008Direct}及 Lawrence C. Evans 的《Partial Differential Equations》\cite{Evans2010Partial}。

\ProseParagraphBreak

极小化并不能找到所有平衡态。一个能量可能沿某些方向趋于负无穷，却在离开原点时先经过一道正的能量屏障。连接原点与低能量点的路径必须越过这道屏障；让所有路径上的最高能量尽可能低，便产生山路构造。\secref{b2:sec:1405}研究这个几何想法何时能产生真正的临界点。

% !TeX root = ../../Book2.tex
\section{泛函与变分}
\label{b2:sec:1401}

固定两端的曲线可以改变形状，而端点条件保持不变。若用函数 \(u\) 表示曲线，一次允许的微小改变便写成 \(u+th\)，其中 \(h\) 在端点为零。长度、面积和能量都是作用于整条曲线的函数；考察它们随 \(t\) 的变化，就把极值问题带回了一元微积分。不过，在函数空间中，沿每个方向分别求导与对所有小扰动作统一估计并不是同一件事。本节先辨明这一差别，再说明积分泛函的变分怎样计算。

\subsection{Gâteaux 导数}
\label{b2:subsec:140101}

设 \(X\) 是实 Banach 空间，\(U\subset X\) 是开集，\(J:U\to\mathbb R\)。给定 \(u\in U\) 和 \(h\in X\)，若极限存在，记
\[
 \delta J(u;h)=\lim_{t\to0}\frac{J(u+th)-J(u)}{t}.
\]
这里的参数 \(t\) 取实数；当未知函数取复值而泛函取实值时，也应把函数空间看作实线性空间。比如 \(J(u)=\lVert u\rVert^2\) 的变分含有内积的实部，并不是复线性泛函。

\Needspace{6\baselineskip}
\begin{definition}[Gâteaux 可微]
\label{b2:def:14-gateaux}
设 \(X\) 是实 Banach 空间，\(J:U\to\mathbb R\)，\(U\subset X\) 开、\(u\in U\)；\(\delta J(u;h)=\lim_{t\to0}(J(u+th)-J(u))/t\)。若各方向的变分都存在，且存在连续线性泛函 \(A\in X^*\)，使得 \(\delta J(u;h)=Ah\) 对所有 \(h\in X\) 成立，则称 \(J\) 在 \(u\) 处 \term[gateaux-kewei]{Gâteaux 可微}[Gâteaux differentiable]，并称 \(A\) 为其 \term[gateaux-daoshu]{Gâteaux 导数}[Gâteaux derivative]，记作 \(D_GJ(u)\)。
\end{definition}

定义中的线性性不是由方向极限自动保证的。在 \(\mathbb R^2\) 上令
\[
 F(x,y)=\begin{cases}x^3/(x^2+y^2),&(x,y)\ne(0,0),\\0,&(x,y)=(0,0).\end{cases}
\]
沿方向 \((a,b)\ne(0,0)\) 的导数是 \(a^3/(a^2+b^2)\)。这个表达式关于方向不是线性的，故 \(F\) 虽有所有方向导数，却不满足这里的 Gâteaux 可微定义。不同文献有时把“方向导数均存在”也称为 Gâteaux 可微；使用结论时，应核对是否同时要求导数为连续线性映射。

\Needspace{6\baselineskip}
\begin{proposition}[极小点的必要条件]
\label{b2:prop:14-gateaux-minimum}
设 \(X\) 是实 Banach 空间，\(V\subset X\) 是线性子空间，\(J\) 是定义在 \(u\in X\) 某开邻域内的实值泛函。若 \(u\) 是 \(J\) 在仿射空间 \(u+V\) 上的局部极小点，且 \(\delta J(u;h)\) 对每个 \(h\in V\) 都存在，则
\[\delta J(u;h)=0\qquad(h\in V).\]
\end{proposition}
\begin{proof}
固定 \(h\in V\)。当 \(t\) 的绝对值充分小时，\(J(u+th)-J(u)\ge0\)。除以正的 \(t\) 再令 \(t\downarrow0\)，得变分非负；除以负的 \(t\) 再令 \(t\uparrow0\)，得变分非正。两个单侧极限相同，故变分为零。
\end{proof}

此处真正起作用的是正负两个方向都可行。若要求 \(u\ge0\)，在接触零值的区域只能作某些单侧扰动，必要条件一般变成不等式。类似地，固定端点要求扰动在端点消失；固定积分则要求扰动的积分为零。求第一变分之前，先写清允许的扰动，比形式地求导更为重要。

\subsection{Fréchet 导数}
\label{b2:subsec:140102}

第一章在有限维空间中讨论过 Fréchet 导数；同一定义适用于 Banach 空间。设 \(X\) 为实 Banach 空间，\(U\subset X\) 开，\(J:U\to\mathbb R\)、\(u\in U\)。泛函 \(J\) 在 \(u\) 处 Fréchet 可微，意味着存在 \(A\in X^*\)，使
\[
 J(u+h)=J(u)+Ah+r(h),\qquad
 \frac{|r(h)|}{\lVert h\rVert_X}\longrightarrow0\quad(h\to0).
\]
这里的余项估计对所有小方向统一成立。因此 Fréchet 可微必然蕴含 Gâteaux 可微，且两种导数相同；反过来则需要额外条件。

\Needspace{6\baselineskip}
\begin{proposition}[导数连续时的统一逼近]
\label{b2:prop:14-continuous-gateaux}
设 \(X\) 是实 Banach 空间，实值泛函 \(J\) 在 \(u\) 的一个开邻域内 Gâteaux 可微。若 \(v\mapsto D_GJ(v)\) 在 \(u\) 处关于 \(X^*\) 的算子范数连续，则 \(J\) 在 \(u\) 处 Fréchet 可微，且 \(DJ(u)=D_GJ(u)\)。
\end{proposition}
\begin{proof}
给定 \(\varepsilon>0\)，选取 \(r>0\)，使 \(\lVert v-u\rVert_X<r\) 时有 \(\lVert D_GJ(v)-D_GJ(u)\rVert_{X^*}<\varepsilon\)。对 \(\lVert h\rVert_X<r\)，在 \([0,1]\) 上考察
\[g(t)=J(u+th)-J(u)-tD_GJ(u)h.\]
它在区间上连续、内部可微，且 \(|g'(t)|\le\varepsilon\lVert h\rVert_X\)。一元中值定理给出 \(|g(1)|\le\varepsilon\lVert h\rVert_X\)，正是所需的统一余项估计。
\end{proof}

即便 Gâteaux 导数在一点是零映射，原函数也未必在该点连续。\subsecref{b2:subsec:010101}已经计算过的反例是
\[
 J(x,y)=\begin{cases}\displaystyle\frac{x^3y}{x^6+y^2},&(x,y)\ne(0,0),\\0,&(x,y)=(0,0).\end{cases}
\]
它在原点的所有方向导数均为零，即 \(D_GJ(0)=0\)，但沿 \(y=x^3\)、\(x\ne0\) 有 \(J(x,x^3)=1/2\)，故原点处不连续，也不可能 Fréchet 可微。这里需要保留的新区别是：每个固定方向的极限存在，并不提供对小扰动统一的余项估计；前一命题增加的导数算子范数连续性，才足以把这些逐方向极限接成同一个局部线性逼近。

\subsection{第一变分}
\label{b2:subsec:140103}

对实 Banach 空间 \(X\) 的开集 \(U\) 上的 Fréchet 可微泛函 \(J:U\to\mathbb R\)，\(u\in U\)、\(h\in X\) 时的 \(DJ(u)h\) 即泛函对扰动的一阶线性响应，称为 \term[diyi-bianfen]{第一变分}[first variation]。对积分泛函，常见计算是“把导数移入积分号”。这一计算需要一致的支配或余项估计；下面先在 \(C^1\) 空间内证明一个便于直接使用的版本。

\Needspace{6\baselineskip}
\begin{proposition}[光滑积分泛函的第一变分]
\label{b2:prop:14-smooth-integral-variation}
设整数 \(d\geq1\)，\(\Omega\subset\mathbb R^d\) 是有界开集。实函数 \(F(x,z,p)\) 以及对 \((z,p)\in\mathbb R\times\mathbb R^d\) 的一阶偏导数在 \(\overline\Omega\times\mathbb R\times\mathbb R^d\) 上连续。赋予实函数空间 \(C^1(\overline\Omega;\mathbb R)\) 以函数及一阶导数的一致范数。则
\[J(u)=\int_\Omega F(x,u(x),\nabla u(x))\,\mathrm dx\]
为 Fréchet 可微泛函，且
\[
 DJ(u)h=\int_\Omega\bigl(F_z(x,u,\nabla u)h+
             F_p(x,u,\nabla u)\cdot\nabla h\bigr)\,\mathrm dx.
\]
其导数关于 \(u\) 在算子范数下连续。
\end{proposition}
\begin{proof}
固定 \(u\)。当 \(\lVert h\rVert_{C^1}\le1\) 且 \(0\le t\le1\) 时，\((x,u(x)+th(x),\nabla u(x)+t\nabla h(x))\) 落在一个固定紧集内。在这个紧集上，\((z,p)\) 的一阶偏导数一致连续；以 \(\omega(s)\) 表示相应连续模，其中 \(\omega(s)\to0\)。沿 \((h,\nabla h)\) 应用一元微积分基本定理，得
\[
 \begin{split}
 &F(x,u+h,\nabla u+\nabla h)-F(x,u,\nabla u)\\
 &\quad-F_z(x,u,\nabla u)h-F_p(x,u,\nabla u)\cdot\nabla h\\
 &=\int_0^1\bigl[D_{(z,p)}F(x,u+th,\nabla u+t\nabla h)
              -D_{(z,p)}F(x,u,\nabla u)\bigr]\cdot(h,\nabla h)\,\mathrm dt.
 \end{split}
\]
积分后，余项的绝对值不超过
\(C_d|\Omega|\cdot\omega(C_d\lVert h\rVert_{C^1})\cdot\lVert h\rVert_{C^1}\)。所写的一阶项是有界线性泛函，故它就是 Fréchet 导数。用同一个紧集上的连续模比较 \(DJ(u+k)\) 与 \(DJ(u)\)，再对 \(\lVert h\rVert_{C^1}\le1\) 取上确界，就得到导数的算子范数连续性。
\end{proof}

例如，在有界开集 \(\Omega\) 上取实值 \(f\in C(\overline\Omega)\)，则 \(F(x,z,p)=\tfrac12|p|^2-f(x)z\) 满足前一命题的连续性条件，在 \(C^1(\overline\Omega;\mathbb R)\) 上给出
\[DJ(u)h=\int_\Omega\nabla u\cdot\nabla h\,\mathrm dx-\int_\Omega fh\,\mathrm dx.\]
改在 \(H^1(\Omega)\) 上讨论时，上一命题不能原样使用：\(H^1\) 范数一般不控制函数的一致范数，非线性项的增长速度必须与 Sobolev 嵌入相容。特别是，写下 \(\int |u|^q\) 之前，须先确认每个允许的 \(u\) 确有这一可积性。

二次能量可以避开一致范数，直接在其自然空间中计算。设 \(\Omega\subset\mathbb R^d\) 为开集，只要求实值 \(f\in L^2(\Omega)\)，令
\[
 J_f(u)=\frac12\int_\Omega|\nabla u|^2\,\mathrm{d}x
          -\int_\Omega fu\,\mathrm{d}x,
 \qquad u\in H^1(\Omega;\mathbb R).
\]
对每个 \(u,h\in H^1(\Omega;\mathbb R)\)，候选线性响应
\[
 A_uh=\int_\Omega\nabla u\cdot\nabla h\,\mathrm{d}x
             -\int_\Omega fh\,\mathrm{d}x
\]
满足 \(|A_uh|\leq(\|\nabla u\|_2+\|f\|_2)\|h\|_{H^1}\)。直接展开平方还给出精确余项
\[
 J_f(u+h)-J_f(u)-A_uh=\frac12\|\nabla h\|_2^2
                       \leq\frac12\|h\|_{H^1}^2.
\]
除以 \(\|h\|_{H^1}\) 后趋零，故 \(DJ_f(u)=A_u\)。而 \(\|A_u-A_v\|_{(H^1)^*}\leq\|\nabla(u-v)\|_2\)，导数在算子范数下连续。这里的 \(f\) 不必连续；\(L^2\) 配对的有界性及二次余项已经给出所需的可微性。本章后面将以这个空间中的能量进行极小化。

梯度幂能量可以直接在自然的可积空间中微分。这里不能用逐点求导代替范数余项估计，因为扰动可能在很小的集合上很大。

\Needspace{6\baselineskip}
\begin{proposition}[梯度幂能量的连续微分]
\label{b2:prop:14-power-frechet}
设整数 \(d\geq1\)，\(\Omega\subset\mathbb R^d\) 为开集。若 \(1<p<\infty\)，则 \(P(u)=p^{-1}\int_\Omega|\nabla u|^p\) 在实空间 \(W^{1,p}(\Omega)\) 上连续 Fréchet 可微，且
\[DP(u)h=\int_\Omega |\nabla u|^{p-2}\nabla u\cdot\nabla h.\]
\end{proposition}
\begin{proof}
置 \(A_p(z)=|z|^{p-2}z\)，并规定 \(A_p(0)=0\)；\(p=2\) 时直接有 \(A_2(z)=z\)。先核验使积分余项一致变小的两个向量估计：
\[
 |A_p(a+b)-A_p(a)|\le
 \begin{cases}
 C_p(|a|+|b|)^{p-2}|b|,&p\ge2,\\
 C_p|b|^{p-1},&1<p<2.
 \end{cases}
\]
在非零点，\(DA_p(z)=|z|^{p-2}I+(p-2)|z|^{p-4}z\otimes z\)，其算子范数不超过 \(C_p|z|^{p-2}\)。当 \(p\ge2\) 时，沿 \(a+tb\) 积分即得第一式，经过零点也可分成两段取极限。当 \(p<2\)、\(|b|\ge |a|/2\) 时，分别估计两端的 \(|A_p|\)，得到第二式；当 \(|b|<|a|/2\) 时，整条线段离零点至少 \(|a|/2\)，导数估计给出 \(C_p|a|^{p-2}|b|\le C_p|b|^{p-1}\)。

对 \(v,w\in L^p(\Omega;\mathbb R^d)\)，Hölder 不等式于是给出
\[
 \|A_p(v+w)-A_p(v)\|_{p'}\le
 \begin{cases}
 C_p(\|v\|_p+\|w\|_p)^{p-2}\|w\|_p,&p\ge2,\\
 C_p\|w\|_p^{p-1},&1<p<2,
 \end{cases}
 \qquad p'=\frac p{p-1}.
\]
第一种情形用的指数关系是 \(1/p'=(p-2)/p+1/p\)；第二种情形直接把 \(|w|^{p-1}\) 的 \(p'\) 次方积分。另有 \(\|A_p(v)\|_{p'}=\|v\|_p^{p-1}\)，故候选导数确为有界线性泛函。

沿逐点线段积分，再用 Tonelli 与 Hölder，余项满足
\[
 |P(u+h)-P(u)-DP(u)h|
 \le \|\nabla h\|_p\sup_{0\le t\le1}
 \|A_p(\nabla u+t\nabla h)-A_p(\nabla u)\|_{p'}.
\]
上界在 \(p\ge2\) 时为 \(O((\|\nabla u\|_p+\|\nabla h\|_p)^{p-2}\|\nabla h\|_p^2)\)，在 \(p<2\) 时为 \(O(\|\nabla h\|_p^p)\)，除以 \(\|h\|_{W^{1,p}}\) 后都趋零。同一向量估计又说明 \(DP(u+h)\to DP(u)\) 于对偶算子范数，完成证明。
\end{proof}

这也辨明了第一与第二变分的差别：\(1<p<2\) 时一阶导数连续，但 \(|z|^p/p\) 的二阶导数在 \(z=0\) 附近发散；不能把二次能量的有界 Hessian 论证照搬到所有幂能量。下面的存在性只需一阶变分和凸性，不需这种额外二阶正则性。

\subsection{第二变分}
\label{b2:subsec:140104}

设 \(X\) 为实 Banach 空间，\(U\subset X\) 开，\(J:U\to\mathbb R\)，\(u\in U\)、\(h\in X\)，且 \(J\) 在 \(u\) 的邻域可微。当导数映射本身在 \(u\) 处 Fréchet 可微时，其导数 \(D^2J(u)\) 是有界双线性型。对同一方向两次求导所得的
\[\delta^2J(u;h)=D^2J(u)[h,h]\]
称为 \term[dier-bianfen]{第二变分}[second variation]。若 \(J\) 在 \(u\) 附近二阶连续 Fréchet 可微，则沿线段积分两次得到
\[
 J(u+h)-J(u)-DJ(u)h
 =\int_0^1(1-t)D^2J(u+th)[h,h]\,\mathrm dt.
\]
二阶项反映驻点附近的弯曲方向；它提供必要条件，也在具有统一正下界时提供充分条件。

\Needspace{6\baselineskip}
\begin{theorem}[第二变分判别]
\label{b2:thm:14-second-variation-test}
设 \(X\) 是实 Banach 空间，\(u\in X\)，\(V\subset X\) 是闭线性子空间，实值泛函 \(J\) 在 \(u\) 的开邻域内二阶连续 Fréchet 可微。若 \(u\) 是 \(J\) 在 \(u+V\) 上的局部极小点，则
\[DJ(u)h=0,\qquad D^2J(u)[h,h]\ge0\quad(h\in V).\]
反之，若 \(DJ(u)|_V=0\)，且存在 \(c>0\)，使
\[D^2J(u)[h,h]\ge c\lVert h\rVert_X^2\quad(h\in V),\]
则 \(u\) 是 \(J\) 在 \(u+V\) 上的严格局部极小点。
\end{theorem}
\begin{proof}
必要性由一元函数 \(t\mapsto J(u+th)\) 的一、二阶极值条件得到。证明充分性时，利用 \(D^2J\) 的算子范数连续性，取 \(r>0\)，使 \(\lVert v-u\rVert_X<r\) 时 \(\lVert D^2J(v)-D^2J(u)\rVert<c/2\)。对 \(h\in V\)、\(0<\lVert h\rVert_X<r\)，上面的积分余项公式给出
\[
 J(u+h)-J(u)\ge\int_0^1(1-t)\frac c2\lVert h\rVert_X^2\,\mathrm dt
 =\frac c4\lVert h\rVert_X^2>0.
\]
\end{proof}

仅有非负第二变分并不足以判别极小：\(t^4\) 与 \(-t^4\) 在零点的第一、第二导数都为零，却有相反的极值性质。在无限维空间里，即使对每个非零 \(h\) 都有 \(D^2J(u)[h,h]>0\)，也不能直接得到定理要求的统一常数 \(c\)。后面处理 Dirichlet 能量时，正是 Poincaré 不等式把梯度能量转成整个函数空间范数的控制，从而弥补这一差别。

% !TeX root = ../../Book2.tex
\section{Euler–Lagrange 方程}
\label{b2:sec:1402}

第一变分为零是函数空间中的驻点条件。它通常以积分等式出现；把扰动的导数移到系数上，才能读出微分方程。这个步骤也同时产生边界项，因此方程与边界条件不能分开猜测。本节先在光滑函数类中推导，下一章再把积分等式本身作为解的定义。

\subsection{一维问题}
\label{b2:subsec:140201}

设 \(a<b\)，考虑端点固定为 \(u(a)=A\)、\(u(b)=B\) 的实函数，及泛函
\[J(u)=\int_a^b F(x,u(x),u'(x))\,\mathrm dx.\]
允许扰动满足 \(h(a)=h(b)=0\)。如果 \(u\) 为极小点，\cref{b2:prop:14-gateaux-minimum,b2:prop:14-smooth-integral-variation}给出
\[\int_a^b\bigl(F_z(x,u,u')h+F_p(x,u,u')h'\bigr)\,\mathrm dx=0.\]

\Needspace{6\baselineskip}
\begin{theorem}[一维 Euler--Lagrange 方程]
\label{b2:thm:14-euler-lagrange-one}
设 \(a<b\)，实值 \(F\in C^2([a,b]\times\mathbb R^2)\)，实值 \(u\in C^2([a,b])\) 是 \(J(v)=\int_a^bF(x,v,v')\,\mathrm dx\) 在 \(v(a)=u(a)\)、\(v(b)=u(b)\) 的函数类中在 \(C^1\) 范数下的局部极小点。则
\[F_z(x,u,u')-\frac{\mathrm d}{\mathrm dx}F_p(x,u,u')=0\quad(a<x<b).\]
更一般地，对任意 \(u\in C^2([a,b])\)，该方程与所有零端点扰动的第一变分为零等价。
\end{theorem}
\begin{proof}
分部积分把第一变分写成 \(\int_a^b q(x)h(x)\,\mathrm dx\)，其中 \(q=F_z-\frac{\mathrm d}{\mathrm dx}F_p\) 连续。若某点 \(q(x_0)>0\)，连续性保证一个小区间内 \(q>0\)；在其中选非负且不恒零的光滑紧支集扰动，便使积分为正，矛盾。对 \(q(x_0)<0\) 同理。因此 \(q=0\)。反向只须把方程代回分部积分等式，边界项因扰动端点为零而消失。
\end{proof}

曲线的弧长对应 \(F=\sqrt{1+p^2}\)，于是 \(\frac{\mathrm d}{\mathrm dx}(u'/\sqrt{1+(u')^2})=0\)。由于 \(p\mapsto p/\sqrt{1+p^2}\) 严格单调，\(u'\) 为常数，驻曲线只能是连接两端点的线段。它确为全局极小曲线，可以再由 \(p\mapsto\sqrt{1+p^2}\) 的凸性证明：以线段斜率 \(m\) 处的切线作下界，积分后线性修正项为 \(\frac{m}{\sqrt{1+m^2}}\int_a^b(u'-m)=0\)。这里微分方程找出候选者，凸性才完成极小性的判断。

\ProseParagraphBreak

若 \(F\) 不显含 \(x\)，沿任何 \(C^2\) 驻曲线有
\[
 \frac{\mathrm d}{\mathrm dx}(F-u'F_p)
 =u'\left(F_z-\frac{\mathrm d}{\mathrm dx}F_p\right)=0.
\]
这给出一个一阶守恒量，常比二阶方程更易使用。但若原式包含 \(x\)，右端还要加上 \(F_x\)，不能继续断言守恒。即便有守恒量，也须重新代入端点条件；局部的一阶积分并不保证全局可解。

\subsection{多维问题}
\label{b2:subsec:140202}

多维区域中，导数向量 \(\nabla h\) 的分部积分产生散度。于是梯度型能量自然对应散度型方程，这也是弱解理论优先采用散度形式的原因。

\Needspace{6\baselineskip}
\begin{theorem}[多维 Euler--Lagrange 方程]
\label{b2:thm:14-euler-lagrange-many}
设整数 \(d\geq1\)，\(\Omega\subset\mathbb R^d\) 是有界开集，实值 \(F\in C^2(\overline\Omega\times\mathbb R\times\mathbb R^d)\)，实值 \(u\in C^2(\Omega)\cap C^1(\overline\Omega)\)。则 \(J(u)=\int_\Omega F(x,u,\nabla u)\,\mathrm dx\) 对每个 \(h\in C_c^\infty(\Omega)\) 的第一变分为零，当且仅当
\[F_z(x,u,\nabla u)-\operatorname{div}\bigl(F_p(x,u,\nabla u)\bigr)=0\quad\text{于 }\Omega.\]
特别地，任一允许这些紧支集扰动的光滑局部极小点都满足此方程。
\end{theorem}
\begin{proof}
第一变分公式对紧支集扰动成立。逐坐标分部积分，边界项为零，得到连续函数 \(F_z-\operatorname{div}F_p\) 与任意测试函数的乘积积分为零。若此函数在一点非零，在小球内选同号的非负凸起函数便产生矛盾，故函数处处为零。反向由同一分部积分恒等式立即得到。
\end{proof}

例如图像面积泛函 \(\int_\Omega\sqrt{1+|\nabla u|^2}\,\mathrm dx\) 的驻点满足
\[\operatorname{div}\frac{\nabla u}{\sqrt{1+|\nabla u|^2}}=0.\]
当梯度很小时，分母接近 \(1\)，线性化给出 \(\Delta u=0\)。这解释了调和函数为何可以描述小斜率的平衡曲面，同时也说明二次 Dirichlet 能量只是面积能量在小梯度附近的近似：大梯度下两者的增长速度和紧性要求不同。

\ProseParagraphBreak

还有一种常见误读：把 \(F_p(x,u,\nabla u)\) 看成只依赖 \(x\) 的显式系数，而漏掉链式法则。以 \(F=\tfrac12a(u)|\nabla u|^2\) 为例，方程是
\[\tfrac12a'(u)|\nabla u|^2-\operatorname{div}(a(u)\nabla u)=0,\]
若展开散度，等价于 \(a(u)\Delta u+\tfrac12a'(u)|\nabla u|^2=0\)，绝不是单纯的 \(a(u)\Delta u=0\)。在粗糙系数或弱正则性情形，保留散度形式通常更可靠。

\subsection{自然边界条件}
\label{b2:subsec:140203}

边界值若预先固定，扰动在边界为零，边界项便不参与驻点条件；边界值若可以自由变化，边界项必须自己消失，由此得到 \term[ziran-bianjie-tiaojian]{自然边界条件}[natural boundary condition]。设整数 \(d\geq1\)，\(\Omega\subset\mathbb R^d\) 为非空有界 \(C^2\) 开集，不要求连通；\(u\) 为实值且在 \(\overline\Omega\) 的邻域内为 \(C^2\)。设实值密度 \(F\) 在 \(\overline\Omega\times\mathbb R\times\mathbb R^d\) 的邻域内为 \(C^2\)，\(B,B_z\) 在 \(\partial\Omega\times\mathbb R\) 上连续。考虑边界能量
\[
 J(u)=\int_\Omega F(x,u,\nabla u)\,\mathrm dx
       +\int_{\partial\Omega}B(x,u)\,\mathrm dS.
\]
对任意实值 \(C^1\) 扰动 \(h\)（定义在闭区域的邻域内），散度公式给出
\[
 DJ(u)h=\int_\Omega(F_z-\operatorname{div}F_p)h\,\mathrm dx
       +\int_{\partial\Omega}(F_p\cdot\nu+B_z)h\,\mathrm dS.
\]
先取紧支集扰动得到内部方程，再允许任意光滑边界迹。这一边界迹自由度可具体核验。把一小块边界写成 \(x_d=\gamma(x')\)，取支集紧含于坐标块的边界函数 \(\psi(x')\)，在邻域中定义
\[h(x',x_d)=\psi(x')\eta(x_d-\gamma(x')),\qquad \eta(0)=1,\]
其中 \(\eta\) 是支集足够小的光滑截断。由于 \(\gamma\in C^2\)，此延拓至少为 \(C^1\)，已足以作第一变分；再乘切向截断并在坐标块外置零，不引入拼接跳跃。有限个坐标块的单位分解可延拓全边界上的允许迹。如果连续边界残差在一点为正，就选择该点附近非负且非零的 \(\psi\)，边界积分为正，与驻定矛盾；为负时同理。因此得到
\[F_p(x,u,\nabla u)\cdot\nu+B_z(x,u)=0\quad\text{于 }\partial\Omega.\]
这里 \(\nu\) 始终是外单位法向。

\ProseParagraphBreak

对于二次模型，不必要求密度对 \(x\) 具有上述全部二阶导数。只取实值 \(f\in C(\overline\Omega)\)、\(\beta,g\in C(\partial\Omega)\)，令 \(F=\tfrac12|p|^2-fz\)、\(B=\tfrac12\beta z^2-gz\)。直接展开二次能量并分部积分，就得到光滑驻点的必要条件
\[-\Delta u=f\quad\text{于 }\Omega,\qquad \partial_\nu u+\beta u=g\quad\text{于 }\partial\Omega.\]
当 \(\beta=0\) 时，这是 Neumann 条件。特别地，完全自由的纯 Dirichlet 能量给出零法向导数，而不是零边界值。若只在一部分边界固定函数值，则在该部分扰动为零，在其余部分产生自然条件；交界处的正则性是另一问题，形式推导不自动保证解在那里光滑。

在纯 Neumann 情形，还必须核查数据相容性。设 \(\Omega_j\) 为 \(\Omega\) 的各连通分支，\cref{b2:prop:13-neumann-compatibility}给出逐分支的必要条件
\[
 \int_{\Omega_j}f\,\mathrm{d}x
       +\int_{\partial\Omega_j}g\,\mathrm{d}S=0.
\]
也可以直接从变分看到这个约束：将 \(u\) 在 \(\Omega_j\) 上加一个常数 \(c\)，梯度能量不变，而
\[
 J(u+c\mathbf1_{\Omega_j})-J(u)
 =-c\left(\int_{\Omega_j}f\,\mathrm{d}x
                 +\int_{\partial\Omega_j}g\,\mathrm{d}S\right).
\]
这些逐分支常数方向在当前 \(C^2\) 边界的光滑变分类中可行，在 \(H^1\) 弱变分类中也可行。因此若有一项不为零，泛函沿该方向就没有驻点，且向一个方向无下界；若全部相容，能量反而对每个分支上的常数平移不变。连通区域须固定一个常数，非连通区域须分别处理各分支的常数自由度。自然边界方程描述的是驻点必须满足的条件，不能代替相容性、极小元存在性和归一化后的唯一性证明。

\ProseParagraphBreak

一维的符号尤其值得核对。对 \(\int_a^bF(x,u,u')\,\mathrm dx+B_a(u(a))+B_b(u(b))\)，两端自由时条件分别为
\[-F_p(a,u(a),u'(a))+B_a'(u(a))=0,\qquad F_p(b,u(b),u'(b))+B_b'(u(b))=0.\]
左端的负号来自外法向朝左。若遗漏它，带弹簧或外力的边界模型就会产生相反的力平衡。

\subsection{向量值问题}
\label{b2:subsec:140204}

现在固定整数 \(d,m\geq1\)、有界开集 \(\Omega\subset\mathbb R^d\)，令 \(u=(u^1,\ldots,u^m)\in C^2(\Omega;\mathbb R^m)\cap C^1(\overline\Omega;\mathbb R^m)\)，把梯度视为矩阵 \(P=(p_i^\alpha)\)，其中 \(p_i^\alpha=\partial_i u^\alpha\)。对实值密度 \(F\in C^2(\overline\Omega\times\mathbb R^m\times\mathbb R^{m\times d})\)，紧支集光滑向量扰动的第一变分为
\[
 DJ(u)h=\int_\Omega\sum_{\alpha=1}^m
 \left(F_{z^\alpha}h^\alpha+\sum_{i=1}^dF_{p_i^\alpha}\partial_i h^\alpha\right)\,\mathrm dx.
\]
各分量扰动可以独立选择。因此，若 \(u\) 光滑且对所有紧支集向量扰动驻定，则逐分量的基本引理给出系统
\[
 F_{z^\alpha}(x,u,Du)-\sum_{i=1}^d\partial_iF_{p_i^\alpha}(x,u,Du)=0,
 \qquad \alpha=1,\ldots,m.
\]
推导只使用上一证明的逐分量分部积分，无须另加各分量彼此独立的假设；它们的耦合已经包含在密度 \(F\) 中。

\ProseParagraphBreak

例如 \(F=\tfrac12\sum_\alpha|\nabla u^\alpha|^2+W(u)\) 给出 \(-\Delta u^\alpha+\partial_\alpha W(u)=0\)。若 \(W(z)=\tfrac12z^{\mathsf T}Mz\)，其中 \(M\) 为实对称矩阵，则为 \(-\Delta u+Mu=0\)。若写入一个非对称矩阵 \(M\)，二次型 \(z^{\mathsf T}Mz\) 实际只保留对称部分；相应变分绝不产生整个非对称矩阵。这一简单计算已经预示：一般线性椭圆方程虽可用双线性型求解，却未必是某个标量二次能量的极值方程。

\ProseParagraphBreak

还须区别值域约束。例如要求 \(|u|=1\) 时，可行的一阶扰动必须满足 \(u\cdot h=0\)。此时只能推出 Euler--Lagrange 向量与所有切向扰动正交，故它沿法向 \(u\)；不能把每个分量直接置零。对光滑的球面约束，乘子实际上可以算出。设整数 \(d,m\geq1\)，\(\Omega\subset\mathbb R^d\) 为非空有界开集，\(u\in C^2(\Omega;\mathbb R^m)\cap C^1(\overline\Omega;\mathbb R^m)\)、\(|u|=1\)，考虑半 Dirichlet 能量。对任意紧支集 \(h\) 且 \(u\cdot h=0\)，使用真正留在球面上的曲线
\[u_t=\frac{u+th}{|u+th|}=\frac{u+th}{\sqrt{1+t^2|h|^2}}.\]
它保持支集外的边界数据，\(\partial_tu_t|_{t=0}=h\)，故驻定给出 \(\int Du:Dh=0\)。对任意紧支集向量 \(\varphi\)，可取 \(h=\varphi-(u\cdot\varphi)u\)。分部积分后说明 \(-\Delta u\) 的切向投影为零，即 \(\Delta u=(u\cdot\Delta u)u\)。另一方面对 \(|u|^2=1\) 微分两次得到
\[0=\tfrac12\Delta|u|^2=|Du|^2+u\cdot\Delta u,\]
于是约束方程为 \(\Delta u+|Du|^2u=0\)。在一维，\(u(x)=(\cos kx,\sin kx)\) 满足 \(u''=-k^2u\)、\(|u'|^2=k^2\)，因而是该约束问题的驻曲线；它通常不是无约束 Dirichlet 问题的驻点。切向扰动确定了光滑驻曲线的约束方程；其极小性还取决于能量在其他可行曲线上的比较。

把这些约束并列，便能看清第一变分为何有时等于零、有时只得到一个方向的不等式。\Cref{b2:tab:14-admissible-variations}中的“边界迹自由”保持区域不动，只让函数的边界值变化；它不是移动区域形状的自由边界问题。约束也可叠加：例如固定迹且固定总积分时，须同时要求 \(h\in H_0^1(\Omega)\) 和 \(\int_\Omega h=0\)。

\begin{table}[htbp]
\centering
\small
\setlength{\tabcolsep}{5pt}
\caption[不同约束下的第一变分方向]{不同约束下的第一变分方向。固定非空有界 Lipschitz 开集 \(\Omega\subset\mathbb R^d\)。除球面行外，均在实空间 \(H^1(\Omega)\) 中讨论，固定迹的参照函数 \(g\in H^1(\Omega)\)，\(M\in\mathbb R\)，\(K\subset H^1(\Omega)\) 为非空闭凸集，泛函 \(J\) 在相应函数的开邻域 Fréchet 可微；经典方程仍按本节光滑积分模型理解。球面行另取 \(C^2\) 向量值积分模型、\(u\in C^2(\Omega;\mathbb R^m)\cap C^1(\overline\Omega;\mathbb R^m)\) 及紧支集 \(C^1\) 扰动，约束在每个空间点成立。各行均列必要条件。}
\label{b2:tab:14-admissible-variations}
\begin{tabularx}{\linewidth}{@{}>{\raggedright\arraybackslash}p{3.1cm}>{\raggedright\arraybackslash}X>{\raggedright\arraybackslash}p{5.7cm}@{}}
\toprule
约束 & 可行的一阶方向 & 第一变分条件 \\
\midrule
固定迹：\(u\in g+H_0^1(\Omega)\) & \(h\in H_0^1(\Omega)\)，\(u+th\) 对正负小 \(t\) 均可行 & \(DJ(u)h=0\) 对所有零迹方向成立 \\
固定积分：\(\int_\Omega u=M\) & \(h\in H^1(\Omega)\)、\(\int_\Omega h=0\)；如另固定迹，还须 \(h\in H_0^1\) & 在零积分子空间上 \(DJ(u)h=0\)，方程可留下常数乘子 \\
边界迹自由，区域固定 & \(h\in H^1(\Omega)\)，边界迹可变化 & 所有方向的第一变分为零；光滑积分模型产生自然边界条件 \\
逐点球面值：\(|u(x)|=1\) & \(u(x)\cdot h(x)=0\)；\(h\) 是可行曲线的速度 & Euler--Lagrange 向量的切向分量为零，不要求整个向量为零 \\
闭凸障碍：\(u\in K\) & \(v-u\)、\(v\in K\)，线段 \(u+t(v-u)\) 在 \(0\leq t\leq1\) 可行 & 极小点满足 \(DJ(u)(v-u)\geq0\)，这是单侧变分不等式 \\
\bottomrule
\end{tabularx}
\end{table}

固定积分行的常数乘子来自一个简单的线性事实。记 \(L(h)=\int_\Omega h\)，在所讨论的线性扰动空间内若 \(L\) 非零，可取 \(h_0\) 使 \(L(h_0)=1\)。当 \(DJ(u)\) 在 \(\ker L\) 上为零时，把 \(h\) 分成 \(h-L(h)h_0\) 与 \(L(h)h_0\)，便有 \(DJ(u)h=\lambda L(h)\)，其中 \(\lambda=DJ(u)h_0\)。

球面行的切向条件只限制曲线的初始速度；一般有 \(|u+th|^2=1+t^2|h|^2\)，所以这条直线不会留在球面上。前面使用的 \((u+th)/|u+th|\) 才是具有该速度的真实可行曲线。闭凸障碍则恰好相反：从 \(u\) 指向每个 \(v\in K\) 的单侧线段由凸性保证可行，除以正参数后得到不等式；只有另加 \(J\) 的凸性时，这组不等式才同时足以保证全局极小。

% !TeX root = ../../Book2.tex
\section{能量泛函}
\label{b2:sec:1403}

同一个偏微分方程可以通过能量的驻点条件出现，但“能量”一词并不总表示待极小化的非负量。椭圆方程常对应空间能量的极小值；波动方程对应时空作用量的驻点，同时另有随时间守恒的正能量。本节用这两个模型把几种不同作用区分开。

\subsection{Dirichlet 能量}
\label{b2:subsec:140301}

设 \(\Omega\subset\mathbb R^d\) 是有界开集，记 \(H_0^1(\Omega)=W_0^{1,2}(\Omega)\)。第一卷定义 22.2.1 将它定义为 \(C_c^\infty(\Omega)\) 在 \(H^1\) 中的闭包，这一定义不要求边界光滑。由第一卷定理 24.1.2 的零边界 Poincaré 不等式，\(\lVert v\rVert_2\le C_P\lVert\nabla v\rVert_2\)，所以 \(\lVert\nabla v\rVert_2\) 是这个空间上的等价 Hilbert 范数。

\ProseParagraphBreak

设 \(\Omega\subset\mathbb R^d\) 是有界开集，给定实值 \(f\in L^2(\Omega)\) 和实值 \(g\in H^1(\Omega)\)，把允许函数写为 \(g+H_0^1(\Omega)\)，考虑 \term[dirichlet-nengliang]{Dirichlet 能量}[Dirichlet energy]及线性外力项组成的泛函
\[
 E(u)=\frac12\int_\Omega|\nabla u|^2\,\mathrm dx-\int_\Omega fu\,\mathrm dx.
\]
仿射空间的写法精确表达“\(u\) 与 \(g\) 具有相同边界数据”：即使边界粗糙到尚未定义逐点边界迹，也可以使用它。

\Needspace{6\baselineskip}
\begin{proposition}[二次能量的驻点与极小点]
\label{b2:prop:14-dirichlet-minimum}
设整数 \(d\geq1\)，\(\Omega\subset\mathbb R^d\) 为非空有界开集，\(f\in L^2(\Omega;\mathbb R)\)、\(g\in H^1(\Omega;\mathbb R)\)，\(E(v)=\frac12\int_\Omega|\nabla v|^2-\int_\Omega fv\)。\(u\in g+H_0^1(\Omega;\mathbb R)\) 是 \(E\) 在该仿射空间上的极小点，当且仅当
\[
 \int_\Omega\nabla u\cdot\nabla h\,\mathrm dx=\int_\Omega fh\,\mathrm dx
 \qquad(h\in H_0^1(\Omega)).
\]
这样的极小点至多有一个，且对每个 \(h\in H_0^1(\Omega)\) 都有
\[E(u+h)-E(u)=\tfrac12\lVert\nabla h\rVert_2^2.\]
\end{proposition}
\begin{proof}
Cauchy--Schwarz 不等式保证所有积分有限。直接展开平方得
\[
 E(u+h)-E(u)=\int_\Omega\nabla u\cdot\nabla h-\int_\Omega fh
              +\tfrac12\lVert\nabla h\rVert_2^2.
\]
其中前两项构成连续线性泛函，最后一项是二阶余项，故一阶导数正是前两项。极小点的双侧变分为零，得到所列积分等式。反之，积分等式使能量差等于非负的梯度平方，所以 \(u\) 为全局极小点。若另一点 \(u+h\) 也极小，则 \(\nabla h=0\)；Poincaré 不等式推出 \(h=0\)。
\end{proof}

这一命题暂时没有断言极小点存在；它说明一旦构造出驻点，就无须再排查鞍点。二次展开还给出误差的精确解释：任何满足同一边界数据的近似函数，其能量超额恰为梯度误差平方的一半。与一般的第二变分局部判别相比，这是全局恒等式，原因在于能量严格二次且对允许方向一致凸。

\ProseParagraphBreak

在区间 \((0,1)\) 上取 \(f=1\)、\(g=0\)。函数 \(u(x)=x(1-x)/2\) 满足 \(-u''=1\) 和零端点条件；对光滑紧支集扰动分部积分，再用密度推广到 \(H_0^1\)，即可验证积分等式。因此它是唯一极小点，能量为 \(-1/24\)。负能量并不矛盾：梯度项非负，外力项却可以降低总泛函值；强制性要求的是大范数下的增长，而不是每一点的非负性。

\subsection{波动作用量与守恒能量}
\label{b2:subsec:140302}

设整数 \(d\geq1\)，\(\Omega\subset\mathbb R^d\) 为非空有界 \(C^1\) 开集，\(c,T>0\)，并考虑在 \([0,T]\times\overline\Omega\) 的开邻域内为 \(C^2\)、在侧边界取零值的实函数 \(u(t,x)\)。时空作用量为
\[
 \mathcal A(u)=\frac12\int_0^T\int_\Omega
       \bigl(|\partial_tu|^2-c^2|\nabla_xu|^2\bigr)\,\mathrm dx\,\mathrm dt.
\]
对在初、末时刻和空间边界均为零的扰动 \(h\)，时、空分别分部积分，得到
\[
 D\mathcal A(u)h=\int_0^T\int_\Omega
       (-\partial_t^2u+c^2\Delta_xu)h\,\mathrm dx\,\mathrm dt.
\]
所以驻点满足波动方程 \(\partial_t^2u-c^2\Delta u=0\)。这里固定初末值是推导变分方程时对扰动的限制；实际波动初值问题仍由初位移和初速度确定，不能把两种问题混为一谈。

\ProseParagraphBreak

作用量的二次型有正负两部分，一般并不在解处取得极小值。例如取乘积扰动 \(h(t,x)=a(t)b(x)\)，其中两因子均光滑紧支集，则
\[
 2\mathcal A(h)=\lVert a'\rVert_2^2\lVert b\rVert_2^2
      -c^2\lVert a\rVert_2^2\lVert\nabla b\rVert_2^2.
\]
固定非零 \(b\) 后，取 \(a_N(t)=\eta(t)\sin(Nt)\)，其中 \(0\ne\eta\in C_c^\infty(0,T)\)。展开正弦平方并用\cref{b2:thm:03-riemann-lebesgue}，有 \(\|a_N\|_2^2\to\|\eta\|_2^2/2\) 及 \(N^{-2}\|a_N'\|_2^2\to\|\eta\|_2^2/2\)，所以作用量为正。反之，固定非零 \(a\)，取 \(b_N(x)=\psi(x)\sin(Nx_1)\)，其中 \(0\ne\psi\in C_c^\infty(\Omega)\)。相同的平方展开给出 \(\|b_N\|_2^2\to\|\psi\|_2^2/2\) 及 \(N^{-2}\|\nabla b_N\|_2^2\to\|\psi\|_2^2/2\)，所以作用量为负。若 \(u\) 为驻点，精确二次展开给出 \(\mathcal A(u+th)-\mathcal A(u)=t^2\mathcal A(h)\)。故可把两类扰动都缩得任意小，在每个驻点附近仍有增加和降低作用量的方向。

\Needspace{6\baselineskip}
\begin{proposition}[波动能量守恒]
\label{b2:prop:14-wave-energy}
设整数 \(d\geq1\)，\(c,T>0\)，\(\Omega\subset\mathbb R^d\) 为非空有界 \(C^1\) 开集，实函数 \(u\) 在 \([0,T]\times\overline\Omega\) 的某开邻域内属于 \(C^2\)，满足 \(u_{tt}-c^2\Delta u=0\) 于 \((0,T)\times\Omega\)，且 \(u(t,y)=0\) 对 \(0\leq t\leq T\)、\(y\in\partial\Omega\) 成立。则
\[
 \mathcal E(t)=\frac12\int_\Omega
          \bigl(|u_t(t,x)|^2+c^2|\nabla u(t,x)|^2\bigr)\,\mathrm dx
\]
与时间无关。
\end{proposition}
\begin{proof}
对时间求导并对空间分部积分，得
\[
 \mathcal E'(t)=\int_\Omega u_t(u_{tt}-c^2\Delta u)\,\mathrm dx
           +c^2\int_{\partial\Omega}u_t\partial_\nu u\,\mathrm dS=0.
\]
内部项由方程为零，边界项由对恒定零边界条件求时间导数所得 \(u_t=0\) 而消失。
\end{proof}

守恒能量中空间项带正号，作用量中带负号。前者衡量一个时刻的状态，后者比较整段时空运动。对同样的区域、实值 \(C^2\) 解及恒定零边界，若方程改为 \(u_{tt}-c^2\Delta u+\gamma u_t=f\)，其中 \(\gamma\geq0\)、\(f\) 连续，则同一计算给出
\[
 \mathcal E'(t)=\int_\Omega f(t,x)u_t(t,x)\,\mathrm{d}x
                    -\gamma\|u_t(t)\|_2^2.
\]
这些等式将外界做功和耗散分开记录。\cref{b2:tab:14-action-energy}集中比较作用量与能量；表中的守恒结论始终包含零边界且无外力、无阻尼的条件。

\begin{table}[htbp]
\centering
\caption[波动作用量与守恒能量的区别]{波动作用量与守恒能量的区别。取整数 \(d\geq1\)、\(c,T>0\)、非空有界 \(C^1\) 开集 \(\Omega\subset\mathbb R^d\)，函数实值且光滑到足以完成正文的微分与分部积分，侧边界为恒定零值。作用量列的变分还固定初、末时刻的函数值；能量列讨论满足波动方程的实际演化。}
\label{b2:tab:14-action-energy}
\small
\begin{tabularx}{\textwidth}{@{}p{5.2em}XX@{}}
\toprule
比较项目 & 时空作用量 \(\mathcal A\) & 守恒能量 \(\mathcal E(t)\)\\
\midrule
积分范围 & 整段 \((0,T)\times\Omega\) & 固定时刻的 \(\Omega\)\\
空间梯度项 & \(-c^2|\nabla u|^2/2\) & \(+c^2|\nabla u|^2/2\)\\
数学作用 & 对允许扰动驻定，推出运动方程 & 控制速度与梯度，进而估计解的差\\
极值或演化 & 驻点一般为鞍点 & 无外力、无阻尼时随时间守恒\\
\bottomrule
\end{tabularx}
\end{table}

\subsection{PDE 与极值问题}
\label{b2:subsec:140303}

从微分方程寻找泛函时，最容易检验的是对称性。设 \(H\) 是实 Hilbert 空间，\(a:H\times H\to\mathbb R\) 是有界双线性型，\(\ell\in H^*\)。令
\[J(u)=\tfrac12a(u,u)-\ell(u).\]
直接展开可得
\[DJ(u)h=\tfrac12\bigl(a(u,h)+a(h,u)\bigr)-\ell(h).\]
因此仅当 \(a\) 对称时，驻点条件才是 \(a(u,h)=\ell(h)\)。反对称部分在 \(a(u,u)\) 中完全消失，却可能实质影响方程的解。

\ProseParagraphBreak

例如 \(H=\mathbb R^2\)，令 \(Au=(u_1-u_2,u_1+u_2)\)，\(a(u,h)=(Au)\cdot h\)。尽管 \(a(u,u)=|u|^2\) 严格为正，极小化 \(\tfrac12|u|^2-f\cdot u\) 得到的是 \(u=f\)，而原方程是 \(Au=f\)。两者一般不同。下一章的 Lax--Milgram 定理直接求解有界强制双线性型，无须先把它变成能量；这正是从变分极小化过渡到一般弱解方法时必须增加的能力。

\ProseParagraphBreak

另一个区别是驻点与动力学。若 \(E(u)=\tfrac12\int|\nabla u|^2\)，形式上的 \(L^2\) 梯度为 \(-\Delta u\)。沿最陡下降运动 \(u_t=-\operatorname{grad}E(u)\)，得到热方程 \(u_t=\Delta u\)，而平衡点满足 Laplace 方程。对光滑零边界解，分部积分给出
\[\frac{\mathrm d}{\mathrm dt}E(u(t))=-\int_\Omega|u_t|^2\,\mathrm dx\le0.\]
这项恒等式说明热演化趋向降低能量，却不在这里自动证明长时间收敛；还需要轨道紧性、极限点的识别等信息。变分结构可以提供估计和候选极限，但不能代替这些额外步骤。

% !TeX root = ../../Book2.tex
\section{直接方法}
\label{b2:sec:1404}

先推导 Euler--Lagrange 方程再求解，有时会把存在性问题变得更难。直接方法反过来考察能量本身：取能量趋近下确界的序列，控制它的大小，抽取极限，再证明极限没有丢失边界条件和能量下界。第一卷定理 9.5.9 已在自反 Banach 空间中证明这一抽象原则。对积分能量而言，关键是把能量有界转为函数空间中的有界，再判断弱极限能否保留约束并降低能量。

\subsection{强制性}
\label{b2:subsec:140401}

在 Banach 空间 \(X\) 的约束集 \(K\) 上，泛函的强制性指
\[J(v)\longrightarrow+\infty\qquad(v\in K,\ \lVert v\rVert_X\to\infty).\]
它保证能量有界的候选者不能逃到范数无穷远处。先取整数 \(d\geq1\)、非空有界开集 \(\Omega\subset\mathbb R^d\) 及实值 \(f\in L^2(\Omega)\)。第一卷定理24.1.2的零边界 Poincaré 不等式说明：只要开集包含于有限宽条带，就有 \(\|v\|_2\leq C_P\|\nabla v\|_2\) 对全部 \(v\in H_0^1(\Omega)\) 成立。有界开集满足这个几何条件，不需要连通性或边界光滑性。因此 Cauchy--Schwarz 不等式给出
\[
 \frac12\lVert\nabla v\rVert_2^2-\int_\Omega fv
 \ge\frac12\lVert\nabla v\rVert_2^2-C_P\lVert f\rVert_2\lVert\nabla v\rVert_2.
\]
右端是开口向上的二次式；又因梯度范数在 \(H_0^1\) 上与完整的 \(H^1\) 范数等价，故这项能量强制。若固定实值 \(g\in H^1(\Omega)\) 并允许边界类 \(g+H_0^1(\Omega)\)，写成 \(u=g+v\)，交叉项 \(\int\nabla g\cdot\nabla v\) 仍只按 \(\lVert\nabla v\rVert_2\) 线性增长，因此同样得到强制性。

\ProseParagraphBreak

边界限制不能从这个论证中删除。在同一有界开集的实空间 \(H^1(\Omega)\) 中，常函数的梯度能量为零，故单靠梯度无法控制范数。若 \(\int_\Omega f\ne0\)，沿适当符号的常函数甚至能让总能量趋向 \(-\infty\)。齐次 Neumann 问题因而至少要求 \(\int_\Omega f=0\)。对有界 \(C^1\) 开集、连续源项及连续边界通量的经典问题，\cref{b2:prop:13-neumann-compatibility}还证明了总源量加总外通量为零的必要条件。相容条件排除能量沿常数方向下降，却还没有固定常数自由度。

\ProseParagraphBreak

对于非空有界连通 Lipschitz 区域 \(\Omega\subset\mathbb R^d\)，第一卷定理24.1.3已证明 Poincaré--Wirtinger 不等式；在本处的 \(p=2\) 情形，它写成
\[
 u_\Omega=\frac1{|\Omega|}\int_\Omega u,\qquad
 \|u-u_\Omega\|_{L^2(\Omega)}
 \leq C_\Omega\|\nabla u\|_{L^2(\Omega)}
 \quad(u\in H^1(\Omega)).
\]
所以在零均值子空间中，梯度可以控制完整范数。连通性在这里不可删除：在 \((0,1)\cup(2,3)\) 上分别取常数 \(1\)、\(-1\)，所得函数总均值为零、弱梯度为零，\(L^2\) 范数却为 \(\sqrt2\)。不连通时应在每个连通分支分别处理相容条件、减去该分支的均值；应用上述估计还要求各分支满足相应的几何条件，若分支无限多，不能未经说明就把各自常数当成统一常数。即使开集连通，也不能只凭有界性调用这一版本。Neumann 零均值估计使用上述 Lipschitz 与连通条件；零边界 \(W_0^{1,p}\) 的存在性则适用于任意非空有界开集。

\subsection{极小化序列}
\label{b2:subsec:140402}

设 \(K\) 为非空集合、\(J:K\to(-\infty,+\infty]\)。若 \(m=\inf_KJ\) 是有限实数，可以按下确界的定义选取 \(u_n\in K\)，满足 \(m\le J(u_n)<m+1/n\)。这样的序列称为极小化序列。选取过程既不要求驻点已存在，也不要求每个 \(u_n\) 近似满足 Euler--Lagrange 方程。

\ProseParagraphBreak

控制能量后，强制性使序列有界，但有界不等于强收敛。无限维 Hilbert 空间中的标准正交序列有界，却没有强收敛子列；它仅弱收敛到零。这就是直接方法引入弱拓扑的实质理由。另一方面，若约束集不是弱闭的，弱极限可能不再可行；若能量不是弱下半连续的，极限能量可能突然升高。序列有界、约束保留与能量比较分别承担不同任务，不能用一句“由紧性”全部略去。

\ProseParagraphBreak

二次能量还有一个可直接看见的特殊现象。取非空有界开集 \(\Omega\subset\mathbb R^d\)、实值 \(f\in L^2(\Omega)\) 与 \(g\in H^1(\Omega)\)，记 \(E(u)=\tfrac12\int_\Omega|\nabla u|^2-\int_\Omega fu\)、\(m=\inf_{g+H_0^1(\Omega)}E\)。若 \(u_n,u_k\in g+H_0^1(\Omega)\)，则
\[
 E\left(\frac{u_n+u_k}{2}\right)
 =\frac{E(u_n)+E(u_k)}2-\frac18\lVert\nabla(u_n-u_k)\rVert_2^2.
\]
左端不小于下确界 \(m\)。因此极小化序列事实上在梯度范数下为 Cauchy 序列，借完备性可以直接得到强极限。一般非二次能量不一定有这个精确恒等式，弱紧性的方法因而更具普遍性。

\subsection{弱紧性}
\label{b2:subsec:140403}

设整数 \(d\geq1\)、\(\Omega\subset\mathbb R^d\) 为非空开集、\(1<p<\infty\)，函数取实值。映射
\[
 T:W^{1,p}(\Omega)\longrightarrow(L^p(\Omega))^{d+1},\qquad
 Tu=(u,\partial_1u,\ldots,\partial_du)
\]
把 Sobolev 空间识别为一个闭子空间。这里要区分空间结构与范数的数值。第一卷采用逐分量范数
\[
 N_p(u)=\left(\|u\|_p^p+\sum_{i=1}^d\|\partial_i u\|_p^p\right)^{1/p};
\]
它正好与乘积空间的 \(\ell^p\) 范数对应。能量计算常使用欧氏梯度，得到另一个范数
\[
 G_p(u)=\left(\|u\|_p^p+\int_\Omega|\nabla u|_2^p\right)^{1/p}.
\]
令 \(\gamma=|1/p-1/2|\)，有限维向量的范数比较逐点给出
\[
 d^{-\gamma}N_p(u)\leq G_p(u)\leq d^\gamma N_p(u).
\]
因此两种范数等价；\(p=2\) 时它们相同，一般 \(p\) 时不能把欧氏梯度能量范数与逐分量乘积范数默认当作等距。本处的自反性论证只需要这种等价性。

\ProseParagraphBreak

像的闭性也包含一个需要保留的导数关系。若 \(Tu_n\to(v_0,v_1,\ldots,v_d)\) 于乘积范数，则对每个 \(\varphi\in C_c^\infty(\Omega)\)，Hölder 不等式允许在
\(\int u_n\partial_i\varphi=-\int(\partial_i u_n)\varphi\) 中取极限，得到
\(\int v_0\partial_i\varphi=-\int v_i\varphi\)。所以 \(v_i=\partial_i v_0\) 是弱导数，极限仍属于像。有限个自反 \(L^p\) 空间的乘积及其闭子空间仍自反，故第一卷定理9.4.7给出：每个有界 \(W^{1,p}(\Omega)\) 序列都有在该空间弱收敛的子列。沿同一子列，连续的分量映射给出
\[
 u_n\rightharpoonup u\ \text{于 }L^p(\Omega),\qquad
 \partial_i u_n\rightharpoonup\partial_i u\ \text{于 }L^p(\Omega)
 \quad(1\leq i\leq d).
\]
反过来，有限个分量的弱收敛也能识别 \(W^{1,p}\) 中的弱收敛：闭子空间上的有界线性泛函可由 Hahn--Banach 定理延拓到乘积空间，再分解为对这些分量的有限次测试。

\ProseParagraphBreak

对非空开集 \(\Omega\subset\mathbb R^d\)、\(1<p<\infty\) 与实值 \(g\in W^{1,p}(\Omega)\)，仿射边界类 \(g+W_0^{1,p}(\Omega)\) 是强闭凸集，因而弱闭。这一弱闭性来自第一卷定理9.5.2的闭凸集结论。若另给实值可测障碍 \(\psi\)，约束 \(u\ge\psi\) 几乎处处也确定一个强闭凸集：强收敛序列可抽取进一步逐点几乎处处收敛的子列，从而保留不等式；闭凸性再给出弱闭性。该障碍类还须非空，才能用来寻找极小元。

\ProseParagraphBreak

端点 \(p=1\) 与 \(p=\infty\) 不在上述自反性结论内。例如在 \((0,1)\) 上，\(u_n(x)=\min\{nx,1\}\) 的 \(W^{1,1}\) 范数有界，但导数 \(n\mathbf1_{(0,1/n)}\) 集中到越来越短的边界层中。若某个子列弱收敛到 \(v\in L^1\)，对任意 \(0<\delta<1\) 用 \(\mathbf1_{(\delta,1)}\) 测试，得到 \(\int_\delta^1v=0\)，从而 \(v=0\) 几乎处处；用常函数 \(1\) 测试却应有 \(\int_0^1v=1\)，矛盾。因此导数没有 \(L^1\) 弱收敛子列。函数本身在内部趋于常数 \(1\)，左端数据已丢失。处理线性增长能量时往往需要有界变差函数和测度导数；不能把 \(p>1\) 的论证直接搬到这个端点。

\subsection{弱下半连续性}
\label{b2:subsec:140404}

第一卷定理9.5.8说明：赋范空间上凸泛函的范数下半连续性等价于弱下半连续性。对积分泛函，可以先证强收敛下的下半连续，再利用凸性改变拓扑。不过在取积分之前，还须确认被积函数可测。

\ProseParagraphBreak

具体地，设整数 \(d\geq1\)、\(\Omega\subset\mathbb R^d\) 为非空开集、\(1\leq p<\infty\)，并设
\[
 F:\Omega\times\mathbb R\times\mathbb R^d\longrightarrow[0,\infty)
\]
对每个固定 \((z,\xi)\)，函数 \(x\mapsto F(x,z,\xi)\) 为 Lebesgue 可测；对每个固定 \(x\)，函数 \((z,\xi)\mapsto F(x,z,\xi)\) 连续且凸。取实值 \(u\in W^{1,p}(\Omega)\)，为 \(u\) 及全部一阶弱导数选定有限值可测代表，在共同的零测集上补为零，记
\(V(x)=(u(x),\partial_1u(x),\ldots,\partial_du(x))\)。将 \(V\) 的每个分量先截断到 \([-n,n]\)，再向下取最近的 \(2^{-n}\) 有理网格点，得到只取有限个有理向量值的简单函数 \(V_n\)。每个取值集合都由可测分量的区间条件构成，故可写成
\[
 V_n=\sum_{q\in\mathcal Q_n}q\,\mathbf1_{A_{n,q}},\qquad
 \mathcal Q_n\subset\mathbb Q^{d+1}\ \text{有限},
\]
其中 \((A_{n,q})_{q\in\mathcal Q_n}\) 是 \(\Omega\) 的可测划分，且 \(V_n(x)\to V(x)\)。于是
\[
 F(x,V_n(x))=\sum_{q\in\mathcal Q_n}
                  \mathbf1_{A_{n,q}}(x)F(x,q)
\]
可测；再由后几个变量的连续性取逐点极限，得到 \(F(x,u(x),\nabla u(x))\) 可测。更换代表只改变一个零测集上的值，故
\[
 I(u)=\int_\Omega F(x,u,\nabla u)\,\mathrm dx\in[0,+\infty]
\]
是良定义的泛函。

\ProseParagraphBreak

下面核验下半连续性。若 \(u_n\to u\) 于 \(W^{1,p}(\Omega)\)，当 \(\liminf I(u_n)=+\infty\) 时不等式自动成立；否则先抽取实现该有限下极限的子列，再从中抽取函数与全部一阶导数同时几乎处处收敛的子列。这里导数只有有限个，可以逐次抽取而使用同一个最终子列。连续性给出被积函数几乎处处收敛，非负性使 Fatou 引理可用，因此
\[
 I(u)\leq\liminf_{n\to\infty}I(u_n).
\]
这证明强下半连续性。对任意 \(0<t<1\)，密度的凸性逐点给出
\(I((1-t)u+tv)\leq(1-t)I(u)+tI(v)\)，所以 \(I\) 凸；第一卷定理9.5.8遂给出弱下半连续性。这一段的下半连续性核验也适用于 \(p=1\)，而取得弱子列的自反性步骤另外要求 \(1<p<\infty\)。

\ProseParagraphBreak

上述密度的非负性可放宽为 \(F(x,z,\xi)\geq-a(x)\)，其中 \(a\in L^1(\Omega)\) 与 \((z,\xi)\) 无关；对 \(F+a\) 应用刚才的证明，再减去有限常数 \(\int a\) 即可。若 \(\ell\) 是同一实 Sobolev 空间上的有界线性泛函，线性项在弱收敛下连续，所以 \(I-\ell\) 仍弱下半连续。这个证明覆盖了 \(\tfrac1p\int|\nabla u|^p\) 及许多非负凸势能，却没有覆盖任意非凸密度。特别是向量值问题中，关于梯度的普通凸性虽为充分条件，却远非最一般的下半连续性条件。

\ProseParagraphBreak

以 \(\ell^2\) 上的 \(J(x)=-\lVert x\rVert^2\) 为例，\(J\) 范数连续，但对标准正交序列 \(e_n\rightharpoonup0\) 有 \(J(0)=0>-1=\lim J(e_n)\)。所以范数连续本身不能代替弱下半连续；前面的凸性在改变拓扑时起了决定作用。

\subsection{极小元的存在性}
\label{b2:subsec:140405}

下面把这些核验汇集到一个非线性模型。它既包含 Dirichlet 能量，也说明直接方法并不依赖显式解或线性叠加。证明的次序见\cref{b2:fig:14-direct-method}：强制性先控制极小化序列，自反性提供弱子列；约束弱闭性使极限可行，弱下半连续性使极限达到下确界。这四项解决存在性。严格凸性在取得极小元之后另行排除第二个极小元，不能代替其中任何一步。

% 本书原创：直接方法的存在性链条与严格凸性的独立职责。
\begin{figure}[htbp]
\centering
\begin{tikzpicture}[>=stealth,line width=.7pt,
 every node/.style={font=\small},
 directstep/.style={draw=AnalysisBlue,fill=AnalysisBlue!6,rounded corners=2pt,
              align=center,text width=6cm,inner sep=4pt,minimum height=1.05cm}]
\node[directstep] (sequence) at (0,0)
 {极小化序列\\\(u_n\in K,\quad J(u_n)\to m\)};
\node[directstep] (bounded) at (0,-1.5)
 {强制性\\\((u_n)\) 在 \(X\) 中有界};
\node[directstep] (weak) at (0,-3)
 {自反性与弱紧性\\\(u_{n_k}\rightharpoonup u\)};
\node[directstep] (limit) at (0,-4.8)
 {约束弱闭与能量弱下半连续\\\(u\in K\)\\\(J(u)\leq\liminf_kJ(u_{n_k})=m\)};
\node[directstep] (exists) at (0,-6.6)
 {极小元存在\\\(J(u)=m\)};
\node[directstep,text width=3.1cm] (unique) at (5.5,-6.6)
 {另加严格凸性\\极小元唯一};
\draw[->] (sequence)--(bounded);
\draw[->] (bounded)--(weak);
\draw[->] (weak)--(limit);
\draw[->] (limit)--(exists);
\draw[->] (exists)--(unique);
\end{tikzpicture}
\caption[直接方法的存在性步骤与唯一性条件]{直接方法的存在性步骤与唯一性条件。图中 \(X\) 为实自反 Banach 空间，\(K\subset X\) 非空且弱闭，\(J:K\to(-\infty,+\infty]\) 强制、弱下半连续且至少在一点有限，\(m=\inf_KJ\in\mathbb R\)。左列得到极小元；右下分支另要求 \(K\) 凸、\(J\) 在其有限能量点之间严格凸，才得到唯一性。}
\label{b2:fig:14-direct-method}
\end{figure}


\Needspace{6\baselineskip}
\begin{theorem}[梯度幂能量的极小元]
\label{b2:thm:14-p-energy-minimum}
设整数 \(d\geq1\)，\(\Omega\subset\mathbb R^d\) 是非空有界开集，\(1<p<\infty\)，\(g\in W^{1,p}(\Omega;\mathbb R)\)，\(f\in L^{p'}(\Omega;\mathbb R)\)，其中 \(p'=p/(p-1)\)。则
\[
 J(u)=\frac1p\int_\Omega|\nabla u|^p\,\mathrm dx-\int_\Omega fu\,\mathrm dx
\]
在 \(g+W_0^{1,p}(\Omega;\mathbb R)\) 上有唯一极小元 \(u\)。这里 \(W_0^{1,p}\) 指 \(C_c^\infty(\Omega;\mathbb R)\) 在 \(W^{1,p}\) 范数下的闭包，不要求边界存在经典迹。极小元满足
\[
 \int_\Omega|\nabla u|^{p-2}\nabla u\cdot\nabla h\,\mathrm dx
     =\int_\Omega fh\,\mathrm dx\qquad(h\in W_0^{1,p}(\Omega;\mathbb R)),
\]
其中在 \(\nabla u=0\) 处约定 \(|\nabla u|^{p-2}\nabla u=0\)。
\end{theorem}
\begin{proof}
令 \(u=g+v\)。第一卷定理24.1.2适用于有界开集上的 \(v\in W_0^{1,p}(\Omega)\)，故 \(\|v\|_p\leq C_P\|\nabla v\|_p\)。由三角不等式，\(\lVert\nabla v\rVert_p\le\lVert\nabla u\rVert_p+\lVert\nabla g\rVert_p\)，且
\[\left|\int_\Omega fu\right|\le\lVert f\rVert_{p'}\bigl(\lVert g\rVert_p+C_P\lVert\nabla v\rVert_p\bigr).\]
记 \(t=\|\nabla u\|_p\)，\(A=C_P\|f\|_{p'}\)、\(B=\|f\|_{p'}(\|g\|_p+C_P\|\nabla g\|_p)\)，则
\[
 J(g+v)\geq\frac1p t^p-At-B.
\]
右端在 \(t\geq0\) 上有有限下界，且因 \(p>1\) 而趋向 \(+\infty\)。若 \(\|v\|_{W^{1,p}}\to\infty\)，Poincaré 不等式先给出 \(\|\nabla v\|_p\to\infty\)，再由三角不等式得到 \(t\to\infty\)。所以 \(J\) 在边界类上强制，而且 \(J(g)\) 有限，下确界 \(m\) 是有限实数。

取极小化序列 \(u_n\in g+W_0^{1,p}(\Omega)\)。由强制性，\((u_n)\) 在 \(W^{1,p}\) 中有界；自反性给出弱子列 \(u_{n_k}\rightharpoonup u\)。可行集是非空强闭凸集，因而弱闭，所以 \(u-g\in W_0^{1,p}\)。上一小节以密度 \(F(x,z,\xi)=|\xi|^p/p\) 核验了梯度幂积分的弱下半连续性，而 \(f\) 的配对是有界线性泛函，故弱连续。于是
\[
 m\leq J(u)\leq\liminf_{k\to\infty}J(u_{n_k})=m.
\]
这就得到极小元，也具体实现了第一卷定理9.5.9的直接法。

再证唯一性。\(\xi\mapsto|\xi|^p\) 严格凸：若两个向量长度不同，三角不等式结合 \(s\mapsto s^p\) 的严格凸性给出中点严格不等式；若长度相同而向量不同，则欧氏三角不等式本身严格。若 \(u_1,u_2\) 都极小，而其梯度在正测度集合上不同，将这一不等式积分，中点的梯度能量就严格更小；外力项是线性的，中点仍可行，故与极小性矛盾。所以梯度几乎处处相等；差属于 \(W_0^{1,p}\)，零边界 Poincaré 不等式推出差为零。

最后固定 \(h\in W_0^{1,p}(\Omega;\mathbb R)\)。有限维函数 \(\xi\mapsto|\xi|^p/p\) 在 \(p>1\) 时为 \(C^1\)，梯度是在原点取零的 \(|\xi|^{p-2}\xi\)。沿线段 \(u+th\) 求导时，\(|t|\le1\) 的导数被
\[C_p\bigl(|\nabla u|^{p-1}+|\nabla h|^{p-1}\bigr)|\nabla h|\]
支配；该函数由 Hölder 不等式可积。支配收敛允许交换求导与积分，得第一变分公式。极小点的正负方向都可行，故第一变分为零，即所列等式。
\end{proof}

当 \(p=2\) 时，这就是 Poisson 方程的弱形式。一般 \(p\) 则对应 \(-\operatorname{div}(|\nabla u|^{p-2}\nabla u)=f\)，通常称为 \(p\)-Laplace 方程。所得极小元首先属于能量空间；若要把弱方程逐点展开，还需进一步取得经典导数。

\ProseParagraphBreak

对非空有界开集 \(\Omega\subset\mathbb R^d\)、\(1<p<\infty\)、实值 \(g\in W^{1,p}(\Omega)\)、\(f\in L^{p'}(\Omega)\)，若再给定连续凸函数 \(\Psi:\mathbb R\to[0,\infty)\)，则可在同一 \(g+W_0^{1,p}\) 类中加入势能 \(\int_\Omega\Psi(u)\)。只要某个可行函数的总能量有限，同样的直接方法仍给出极小元；非负势能保留强制性，前述可测性与 Fatou 论证给出弱下半连续性，严格凸的梯度项保留唯一性。但要写出含 \(\Psi'(u)\) 的弱方程，还须补充可微性和增长条件，确保扰动可积并允许积分下求导。存在性所需的下半连续性通常比推导经典微分方程所需的光滑性弱得多，这也是直接方法的重要优势。


非凸密度的困难可以在区间上看见。令 \(w_n\) 在每个长度 \(1/n\) 的小区间上先以斜率 \(1\) 上升半个区间，再以斜率 \(-1\) 下降，且 \(w_n(0)=0\)。则 \(w_n\in W_0^{1,4}(0,1)\)、\(\|w_n\|_\infty\le1/(2n)\)，而 \(|w_n'|=1\) 几乎处处。对 \(C^1\) 测试函数 \(\psi\)，在各小区间减去左端的 \(\psi\) 值，利用 \(w_n'\) 的零均值，得到
\[\left|\int_0^1w_n'\psi\right|\le\frac{\|\psi'\|_\infty}{n}.\]
\(\|w_n'\|_4=1\) 与 \(C^1\) 在 \(L^{4/3}\) 中的稠密性把结论推广到全部对偶测试，所以 \(w_n\rightharpoonup0\) 于 \(W^{1,4}\)。但密度 \(F(p)=(p^2-1)^2\) 给出
\[\int_0^1F(w_n')=0,\qquad \int_0^1F(0)=1.\]
该泛函因而不弱下半连续。范数有界、边界值保留，甚至函数一致收敛，都不足以保留这项能量；弱极限把快速交替的两种斜率平均成零。先以斜率一上升、再以斜率负一下降的帐篷函数仍能取得零能量，因此这个泛函有极小元；失败的是极小化序列的弱极限保持最小能量这一步。前面凸密度的不等式恰好阻止这种能量降落。

% !TeX root = ../../Book2.tex
\OptionalSection{山路定理与 Palais–Smale 条件}{b2:sec:1405}

直接方法寻找能量最低的状态。然而，一项能量可能沿某些方向趋于负无穷，同时在零点附近先升高；这时全局极小元不存在，方程却仍可能有非零解。一个一维图像是
\[
 J(t)=\frac12t^2-\frac14t^4.
\]
零点是局部极小点，\(J(2)<0\)，从 \(0\) 连续走到 \(2\) 必须经过 \(t=1\)，在那里能量达到 \(1/4\)，导数又恰好为零。\cref{b2:fig:14-mountain-profile}画出了曲线先上升、再下降到零以下的一段。函数空间中可以沿很多不同路径绕行；我们希望把各条路径的最高能量尽量压低，找到仍然无法绕开的临界高度。

\begin{figure}[htbp]
\centering
\begin{tikzpicture}[x=4cm,y=4.2cm,font=\small]
 \draw[->] (-0.08,0)--(1.76,0) node[right]{\(t\)};
 \draw[->] (0,-0.42)--(0,0.40) node[above]{\(J(t)\)};
 \draw[AnalysisBlue,thick,domain=0:1.61,samples=90]
   plot (\x,{0.5*\x*\x-0.25*\x*\x*\x*\x});
 \draw[densely dashed,gray] (0,0.25)--(1,0.25)--(1,0);
 \fill[AnalysisBlue] (1,0.25) circle (1.5pt);
 \node[left] at (0,0.25) {\(\frac14\)};
 \node[below] at (1,0) {\(1\)};
 \node[below left] at (0,0) {\(0\)};
 \node[above right] at (1,0.25) {必经的高度};
 \node[above] at (1.4142,0) {\(\sqrt2\)};
\end{tikzpicture}
\caption[一维能量的山口图像]{一维能量 \(J(t)=t^2/2-t^4/4\) 的一段曲线。从 \(0\) 到大于 \(\sqrt2\) 的正数，任何连续路径都经过 \(t=1\)，因而不能始终低于能量 \(1/4\)。}
\label{b2:fig:14-mountain-profile}
\end{figure}

这个想法有两个困难。首先，越来越接近最优高度的路径未必收敛；其次，路径上的高点也未必有收敛子列。Palais–Smale 条件解决第二个困难，下降流则把“不是临界点”转化为“还能继续降低路径”。以下采用实 Hilbert 空间中梯度局部 Lipschitz 的版本，并在一个次临界幂非线性模型上逐项验证这些条件。

\subsection{Palais–Smale 条件与紧性}
\label{b2:subsec:140501}

设 \(H\) 为实 Hilbert 空间，内积记为 \((\cdot,\cdot)_H\)。对 \(J\in C^1(H;\mathbb R)\)，第一卷定理8.4.7的 Riesz 表示把连续线性泛函 \(DJ(u)\) 写成
\[
 DJ(u)h=(\nabla J(u),h)_H,
 \quad \|DJ(u)\|_{H^*}=\|\nabla J(u)\|_H.
\]
这里的梯度由所选 Hilbert 内积确定。例如，能量空间 \(H_0^1\) 的梯度通常不是逐点计算的偏微分表达式 \(-\Delta u-g(u)\)；后者首先是对偶空间中的对象，要经过 Riesz 对应才成为 \(H_0^1\) 中的向量。\(DJ(u)=0\) 的点仍称为驻点，也称临界点；其函数值称为临界值。

\begin{definition}[Palais–Smale 条件]\label{b2:def:14-palais-smale}
设 \(H\) 为实 Hilbert 空间，\(J\in C^1(H;\mathbb R)\)，\(c\in\mathbb R\)。若每个满足
\[
 J(u_n)\longrightarrow c,
 \quad \|DJ(u_n)\|_{H^*}\longrightarrow0
\]
的序列都具有在 \(H\) 范数中收敛的子列，就称 \(J\) 在高度 \(c\) 满足\term[palais-smale-tiaojian]{Palais–Smale 条件}[Palais--Smale condition]，记为 \((\mathrm{PS})_c\)。若只要求 \((J(u_n))\) 有界、\(\|DJ(u_n)\|_{H^*}\to0\) 就能得到范数收敛子列，则称 \(J\) 满足 Palais–Smale 条件，简称 \((\mathrm{PS})\)。满足前面两个极限的序列称为高度 \(c\) 的 Palais–Smale 序列。
\end{definition}

条件同时记录“高度有限”和“各方向斜率趋零”。它并不要求每个能量有界序列都收敛，也不把弱收敛当成结论。若 \(u_{n_k}\to u\) 于 \(H\)，连续性给出 \(J(u)=c\)、\(DJ(u)=0\)，所以这个紧性条件确实把近似驻点变为驻点。全局 \((\mathrm{PS})\) 等价于每个有限 \(c\) 上的 \((\mathrm{PS})_c\)：一方面直接包含，另一方面先从有界实数列 \(J(u_n)\) 抽出一个收敛子列，再在所得高度应用条件。

\ProseParagraphBreak

无限维中，“序列有界”还不足以验证 \((\mathrm{PS})\)。在实空间 \(\ell^2\) 上令
\[
 J(x)=\frac12\sum_{j=1}^{\infty}\frac{x_j^2}{j}.
\]
标准单位向量 \(e_n\) 满足 \(J(e_n)=1/(2n)\to0\)、\(\|\nabla J(e_n)\|=1/n\to0\)，但不同两项距离为 \(\sqrt2\)。因此 \((\mathrm{PS})_0\) 失败，尽管这组近似驻点始终位于单位球上。后面的椭圆模型将先取得有界性，再借紧嵌入和方程残差得到能量范数中的强收敛。

\begin{lemma}[非临界高度附近的梯度下界]\label{b2:lem:14-ps-gradient-gap}
设 \(H\) 为实 Hilbert 空间，\(J\in C^1(H;\mathbb R)\)，\(c\in\mathbb R\)。假定 \(J\) 满足 \((\mathrm{PS})_c\)，且不存在同时满足 \(J(u)=c\)、\(DJ(u)=0\) 的 \(u\in H\)。则存在 \(a>0\)、\(\delta>0\)，使
\[
 |J(u)-c|\le a\quad\Longrightarrow\quad
 \|\nabla J(u)\|_H\ge\delta.
\]
\end{lemma}
\begin{proof}
若不存在这样的 \(a,\delta\)，则对每个正整数 \(n\) 可以找到 \(u_n\)，使 \(|J(u_n)-c|\le1/n\)、\(\|DJ(u_n)\|_{H^*}<1/n\)。由 \((\mathrm{PS})_c\) 抽出范数收敛子列，极限 \(u\) 满足 \(J(u)=c\)、\(DJ(u)=0\)，与假设矛盾。
\end{proof}

在有限维紧集上，连续且处处为正的梯度长度有正下界；这里的能量带通常不紧。Palais–Smale 条件提供了另一种办法：假如统一下界失效，就会产生一组近似驻点，而它们的极限正是被排除的临界点。

\subsection{变形引理}
\label{b2:subsec:140502}

当 \(\nabla J(u)\ne0\) 时，方向 \(-\nabla J(u)\) 降低能量。为了把所有点在同一段时间内降低相近的高度，我们改用
\[
 -\frac{\nabla J(u)}{\|\nabla J(u)\|_H^2}.
\]
沿这个向量场，能量的时间导数恰为 \(-1\)。分母也说明了为什么需要上一引理：若梯度过小，下降速度在空间中可能失去控制。我们还要在目标高度之外关掉运动，使路径的低能量端点保持原位。

\begin{lemma}[有界局部 Lipschitz 向量场的流]\label{b2:lem:14-hilbert-flow}
设 \(H\) 为实 Hilbert 空间，\(F:H\to H\) 局部 Lipschitz，且存在 \(M<\infty\) 使 \(\|F(u)\|_H\le M\) 对全部 \(u\in H\) 成立。对每个 \(u\in H\)，初值问题
\[
 \dot z(t)=F(z(t)),\quad z(0)=u
\]
在全部 \(t\ge0\) 上有唯一 \(C^1\) 解。记解为 \(\eta(t,u)\)，则 \(\eta:[0,\infty)\times H\to H\) 连续，且 \(\|\eta(t,u)-u\|_H\le Mt\)。若 \(F(u)=0\)，则 \(\eta(t,u)=u\)。
\end{lemma}
\begin{proof}
在初值附近选取半径 \(r>0\) 的闭球，使 \(F\) 在稍大的邻域内有 Lipschitz 常数 \(L\)。对足够小的 \(T>0\)，令 \(MT<r\)、\(LT<1\)。完备空间 \(C([0,T];H)\) 中取闭集
\(\{z:\sup_t\|z(t)-u\|_H\le r\}\)，映射
\[
 (\mathcal Tz)(t)=u+\int_0^tF(z(s))\,\mathrm ds
\]
保持该集，并且 Lipschitz 常数不超过 \(q=LT<1\)。这里连续的 \(H\) 值函数的积分可由 Riemann 和在完备空间中取极限定义。从常值曲线 \(z_0(t)=u\) 开始迭代 \(z_{k+1}=\mathcal Tz_k\)，则
\[
 \|z_{k+1}-z_k\|_\infty\le q^k\|z_1-z_0\|_\infty.
\]
几何级数估计使 \((z_k)\) 成为 Cauchy 列；其极限 \(z\) 由 \(\mathcal T\) 连续性满足 \(z=\mathcal Tz\)。两个不动点的距离不超过自身的 \(q\) 倍，故只能相同。积分基本定理再给出 \(C^1\) 解。两个局部解在共同存在区间相同，因而可以逐段拼接。

若最大正向存在时间为有限的 \(T_*\)，由 \(\|\dot z\|_H\le M\) 有
\(\|z(t)-z(s)\|_H\le M|t-s|\)。于是当 \(t\uparrow T_*\) 时，\(z(t)\) 在 \(H\) 中有极限。从该极限重新解局部初值问题，就把原解延长到 \(T_*\) 以后，矛盾。因此解全局存在，积分同一估计即得位移界。

连续依赖也来自局部 Lipschitz 估计。在一个共同的 Lipschitz 邻域内，两条轨道满足
\[
 \|z(t)-\widetilde z(t)\|_H
 \le \|z(0)-\widetilde z(0)\|_H
      +L\int_0^t\|z(s)-\widetilde z(s)\|_H\,\mathrm ds,
\]
故由\cref{b2:lem:ode-gronwall}不超过 \(e^{Lt}\|z(0)-\widetilde z(0)\|_H\)。该引理用于此处连续非负的距离函数，其积分常数正是 \(L\)。固定轨道在有限时间区间上的像是紧集，可用有限个这样的邻域覆盖，再把时间分成有限段依次使用该估计。初值足够接近时，各段轨道仍停留在所选邻域内，于是得到初值的连续依赖；与 \(\|\dot z\|_H\le M\) 合并，得到时间与初值的联合连续性。最后，若 \(F(u)=0\)，常值轨道已经满足初值问题，唯一性给出结论。
\end{proof}

\begin{lemma}[变形引理]\label{b2:lem:14-deformation}
设 \(H\) 为实 Hilbert 空间，\(J\in C^1(H;\mathbb R)\)，且 \(\nabla J:H\to H\) 局部 Lipschitz。设 \(c\in\mathbb R\)、\(\varepsilon>0\)、\(\delta>0\)，并且
\[
 |J(u)-c|\le3\varepsilon
 \quad\Longrightarrow\quad\|\nabla J(u)\|_H\ge\delta.
\]
则存在连续映射 \(\Phi:H\to H\)，满足
\begin{enumerate}
\item\label{b2:item:14-deform-fixed}
若 \(|J(u)-c|\ge2\varepsilon\)，则 \(\Phi(u)=u\)；
\item\label{b2:item:14-deform-decrease}
对全部 \(u\in H\)，有 \(J(\Phi(u))\le J(u)\)；
\item\label{b2:item:14-deform-band}
若 \(J(u)\le c+\varepsilon\)，则 \(J(\Phi(u))\le c-\varepsilon\)。
\end{enumerate}
\end{lemma}
\begin{proof}
取 Lipschitz 函数 \(\chi:\mathbb R\to[0,1]\)，在 \([c-\varepsilon,c+\varepsilon]\) 上为一，在 \(( -\infty,c-2\varepsilon]\cup[c+2\varepsilon,\infty)\) 上为零，中间两段线性连接。在开带 \(|J(u)-c|<3\varepsilon\) 上令
\[
 F(u)=-\chi(J(u))\frac{\nabla J(u)}{\|\nabla J(u)\|_H^2},
\]
在其余点令 \(F(u)=0\)。开带内分母有正下界，所以 \(F\) 局部 Lipschitz。开带边缘附近的 \(\chi(J(u))\) 恒为零，两种定义相容，故 \(F\) 在整个 \(H\) 上局部 Lipschitz，并有 \(\|F\|_H\le1/\delta\)。

由\cref{b2:lem:14-hilbert-flow}得到流 \(\eta(t,u)\)。链式法则给出
\[
 \frac{\mathrm d}{\mathrm dt}J(\eta(t,u))
 =-\chi(J(\eta(t,u)))\le0.
\]
定义 \(\Phi(u)=\eta(2\varepsilon,u)\)。连续性以及能量不增已经成立；在 \(|J(u)-c|\ge2\varepsilon\) 的点，\(F(u)=0\)，所以这些点固定。

最后设 \(J(u)\le c+\varepsilon\)。若起点已经不高于 \(c-\varepsilon\)，能量不增即可。否则，只要轨道能量还高于 \(c-\varepsilon\)，它就处于 \((c-\varepsilon,c+\varepsilon]\)，在那里导数恒为 \(-1\)。它因此在不超过 \(2\varepsilon\) 的时间内降到 \(c-\varepsilon\)，以后仍不增加。这证明最后一项。
\end{proof}

上述结论称为\term[bianxing-yinli]{变形引理}[deformation lemma]。它改变的是整条路径，而不是只挑出路径的一个最大点；这使变形后的对象仍是一条连续路径。局部 Lipschitz 梯度保证了流的唯一性和对初值的连续依赖。一般 \(C^1\) 泛函的梯度只保证连续，不能直接套用这个流引理；更一般的山路理论用局部 Lipschitz 的伪梯度向量场替代原梯度。Antonio Ambrosetti 的《Critical points and nonlinear variational problems》\cite[Sections 2--3]{Ambrosetti1992Critical}讨论了相应的紧性、变形和极小极大方法。

\subsection{山路定理}
\label{b2:subsec:140503}

设零点周围有一圈能量至少为 \(\alpha>0\) 的球面，而球面外的某个点 \(e\) 能量不高于零。每条连接 \(0\) 与 \(e\) 的路径都必须穿过这圈球面；我们能够改变穿越位置，却不能让整条路径都处在 \(\alpha\) 以下。\term[shanlu-dingli]{山路定理}[mountain pass theorem]把这一几何障碍转化为临界点。

\begin{theorem}[山路定理，局部 Lipschitz 梯度版本]\label{b2:thm:14-mountain-pass}
设 \(H\) 为实 Hilbert 空间，\(J\in C^1(H;\mathbb R)\)，\(\nabla J:H\to H\) 局部 Lipschitz，且 \(J(0)=0\)。假定存在 \(\rho>0\)、\(\alpha>0\)、\(e\in H\)，满足
\[
 J(u)\ge\alpha\quad(\|u\|_H=\rho),
 \quad \|e\|_H>\rho,
 \quad J(e)\le0.
\]
令
\[
 \Gamma=\{\gamma\in C([0,1];H):\gamma(0)=0,\ \gamma(1)=e\},
 \quad c=\inf_{\gamma\in\Gamma}\max_{0\le t\le1}J(\gamma(t)).
\]
若 \(J\) 在这个 \(c\) 上满足 \((\mathrm{PS})_c\)，则存在 \(u_*\in H\) 使
\[
 DJ(u_*)=0,\quad J(u_*)=c\ge\alpha>0.
\]
特别地，\(u_*\ne0\)，且 \(u_*\ne e\)。
\end{theorem}
\begin{proof}
线段 \(\gamma(t)=te\) 属于 \(\Gamma\)，而 \(J\circ\gamma\) 在紧区间上有有限最大值，因此 \(c<\infty\)。任意 \(\gamma\in\Gamma\) 的范数从 \(0\) 连续变到大于 \(\rho\)，介值定理给出某个 \(t\) 使 \(\|\gamma(t)\|_H=\rho\)。故每条路径的最大能量至少为 \(\alpha\)，从而 \(c\ge\alpha>0\)。

反设高度 \(c\) 上没有临界点。\cref{b2:lem:14-ps-gradient-gap}给出能量带宽 \(a>0\) 与梯度下界 \(\delta>0\)。取 \(\varepsilon>0\) 使 \(3\varepsilon<a\)、\(2\varepsilon<c\)，应用\cref{b2:lem:14-deformation}得到 \(\Phi\)。两个端点的能量都不高于零，因而小于 \(c-2\varepsilon\)，所以 \(\Phi(0)=0\)、\(\Phi(e)=e\)。

按下确界的定义，可取 \(\gamma\in\Gamma\) 满足 \(\max_tJ(\gamma(t))<c+\varepsilon\)。复合路径 \(\widetilde\gamma=\Phi\circ\gamma\) 连续且端点不变，因此仍属于 \(\Gamma\)。但变形引理给出
\[
 \max_{0\le t\le1}J(\widetilde\gamma(t))\le c-\varepsilon<c,
\]
与 \(c\) 是全部此类路径最大值的下确界矛盾。于是高度 \(c\) 上有临界点。其能量严格高于两个端点，故与它们不同。
\end{proof}

证明没有要求存在一条“最佳路径”。我们只取最大能量小于 \(c+\varepsilon\) 的近似路径；若高度 \(c\) 处没有驻点，就把整条近似路径向下推，造成矛盾。在开头的一维例子中，取 \(e=2\)，每条路径都经过 \(1\)，线段路径又实现最大值 \(1/4\)，故 \(c=1/4\)。无限维的结论保留这种过山口的含义，但通常不能据此断言所得临界点唯一，也不能只凭名称判定它的二阶导数有几个负方向。

\subsection{次临界半线性椭圆模型}
\label{b2:subsec:140504}

现在考虑
\[
 -\Delta u=|u|^{p-2}u\quad\text{于 }\Omega,
 \quad u\in H_0^1(\Omega;\mathbb R).
\]
这里的右端与线性 Poisson 方程相反，来自能量中的负项：
\[
 J(u)=\frac12\int_\Omega|\nabla u|^2\,\mathrm dx
          -\frac1p\int_\Omega|u|^p\,\mathrm dx.
\]
当 \(p>2\) 时，小振幅的二次梯度能量占优势，大振幅的负 \(p\) 次项占优势。这给出山路几何；严格低于 Sobolev 临界指数则提供紧性。

\begin{theorem}[次临界幂方程的非零弱解]\label{b2:thm:14-subcritical-mountain-model}
设整数 \(d\ge1\)，\(\Omega\subset\mathbb R^d\) 为非空有界开集，并令
\[
 2^*=\begin{cases}2d/(d-2),&d\ge3,\\+\infty,&d=1,2.\end{cases}
\]
若 \(2<p<2^*\)，则存在非零的实函数 \(u\in H_0^1(\Omega)\)，使
\[
 \int_\Omega\nabla u\cdot\nabla v\,\mathrm dx
   =\int_\Omega|u|^{p-2}uv\,\mathrm dx
 \quad(v\in H_0^1(\Omega;\mathbb R)).
\]
零边界条件按 \(H_0^1\) 的闭包定义理解。该解可取为能量
\(J(u)=\tfrac12\|\nabla u\|_2^2-\tfrac1p\|u\|_p^p\)
的一个正能量山路临界点。
\end{theorem}
\begin{proof}
在 \(H=H_0^1(\Omega;\mathbb R)\) 上取
\((u,v)_H=\int_\Omega\nabla u\cdot\nabla v\)。第一卷定理24.1.2的零边界 Poincaré 不等式说明梯度范数与完整 \(H^1\) 范数等价，所以这是 Hilbert 空间。第一卷推论25.2.6还说明：在任意有界开集上，零边界嵌入 \(H\to L^p(\Omega)\) 在当前严格次临界指数下是紧的。特别有
\[
 \|u\|_p\le C\|u\|_H.
\]
这个零边界版本适用于这里的有界开集，无需定义逐点边界值。

先验证泛函及其导数。记 \(N(s)=|s|^{p-2}s\)。由一维微积分，
\[
 |N(a)-N(b)|\le(p-1)(|a|+|b|)^{p-2}|a-b|
 \quad(a,b\in\mathbb R).
\]
对函数使用 Hölder 不等式，得到
\begin{equation}\label{b2:eq:14-power-lipschitz}
 \|N(u)-N(w)\|_{p'}
 \le(p-1)(\|u\|_p+\|w\|_p)^{p-2}\|u-w\|_p,
 \quad p'=\frac p{p-1}.
\end{equation}
例如其中的三个指数满足 \((p-2)/p+1/p=1/p'\)。沿 \(u+th\) 积分一阶导数，\eqref{b2:eq:14-power-lipschitz}使余项在有界球内不超过常数乘 \(\|h\|_p^2\)。因此 \(J\) Fréchet 可微，且
\[
 DJ(u)v=(u,v)_H-\int_\Omega N(u)v,
 \quad
 \|DJ(u)-DJ(w)\|_{H^*}\le C_R\|u-w\|_H
\]
在 \(\|u\|_H,\|w\|_H\le R\) 时成立。Riesz 对应等距，故梯度在每个有界球上 Lipschitz。

其次验证山路几何。嵌入估计给出
\[
 J(u)\ge\frac12\|u\|_H^2-\frac{C^p}{p}\|u\|_H^p.
\]
因 \(p>2\)，可选小的 \(\rho>0\)，使右端在 \(\|u\|_H=\rho\) 时至少为 \(\alpha=\rho^2/4>0\)。再选非零 \(\varphi\in C_c^\infty(\Omega;\mathbb R)\)。沿射线有
\[
 J(t\varphi)=\frac{t^2}2\|\varphi\|_H^2
                  -\frac{t^p}p\|\varphi\|_p^p\longrightarrow-\infty
 \quad(t\to+\infty).
\]
故可取足够大的 \(T\)，使 \(e=T\varphi\) 满足 \(\|e\|_H>\rho\)、\(J(e)<0\)。

最后证明 \((\mathrm{PS})\)。设 \((J(u_n))\) 有界、\(\|DJ(u_n)\|_{H^*}\to0\)。消去非线性项得到关键恒等式
\[
 J(u_n)-\frac1pDJ(u_n)u_n
       =\left(\frac12-\frac1p\right)\|u_n\|_H^2.
\]
因此存在 \(C<\infty\) 与 \(\varepsilon_n\to0\)，使
\[
 \left(\frac12-\frac1p\right)\|u_n\|_H^2
       \le C+\varepsilon_n\|u_n\|_H.
\]
左边的系数为正，二次式估计说明 \((u_n)\) 有界。由第一卷定理9.4.7的自反空间弱子列结论，抽取 \(u_n\rightharpoonup u\) 于 \(H\)；再由上述紧嵌入，沿进一步子列有 \(u_n\to u\) 于 \(L^p\)。强极限与弱极限相同，因为连续嵌入保留弱收敛。

现在以 \(u_n-u\) 代入近似方程。其 \(H\) 范数有界，所以
\(DJ(u_n)(u_n-u)\to0\)。另外，Hölder 不等式给出
\[
 \left|\int_\Omega N(u_n)(u_n-u)\right|
 \le\|u_n\|_p^{p-1}\|u_n-u\|_p\longrightarrow0.
\]
于是 \((u_n,u_n-u)_H\to0\)。弱收敛又给出 \((u,u_n-u)_H\to0\)，两式相减便是
\(\|u_n-u\|_H^2\to0\)。这证明了所需的范数紧性。

应用\cref{b2:thm:14-mountain-pass}得到 \(DJ(u)=0\)、\(J(u)>0\)。前面的导数公式给出全部 \(H_0^1\) 测试的弱等式，正能量排除了 \(u=0\)。
\end{proof}

次临界性在证明中有一个精确的使用位置：把有界弱收敛序列变成 \(L^p\) 强收敛序列，从而控制非线性项。当 \(d\ge3\)、\(p=2^*\) 时，连续嵌入仍然成立，紧嵌入却不成立。取内部点 \(x_0\) 和固定非零 \(\varphi\in C_c^\infty(\mathbb R^d)\)，对足够小的 \(\varepsilon\) 令
\[
 u_\varepsilon(x)=\varepsilon^{-(d-2)/2}
               \varphi\left(\frac{x-x_0}{\varepsilon}\right).
\]
其支集落在 \(\Omega\)，而 \(\|\nabla u_\varepsilon\|_2\)、\(\|u_\varepsilon\|_{2^*}\) 都保持不变，函数却在 \(x\ne x_0\) 处最终为零。这组集中函数没有 \(L^{2^*}\) 强收敛子列。它解释了紧性步骤的障碍；要判断临界方程的存在性，还须结合能量高度、区域几何或另外的紧性机制。

\ProseParagraphBreak

所得结论是自然能量空间中的非零弱解。它没有自动给出符号、唯一性或经典二阶正则性；这些问题分别需要正部截断、比较原理或椭圆估计等进一步分析。对 \(p>2\) 而言，直接极小化当前能量会沿射线降到负无穷，山路方法却从零点附近的正能垒中找到了另一种驻点。这是极小极大方法与极小化方法在这个模型中的实际分工。

% 本章小结（不计入编号）
\begin{chaptersummary}
第一变分把极小问题转成对全部允许方向成立的积分等式；分部积分把该等式拆成内部 Euler--Lagrange 方程和自然边界条件。允许的方向由约束决定，故固定边界、自由边界及单侧约束产生不同形式。第二变分非负只是必要条件，一致正下界才给出局部极小性的稳定判据。对凸二次能量，精确能量差进一步把驻点与全局唯一极小点联系起来。

\ProseParagraphBreak

直接方法从极小化序列出发。强制性控制序列，Sobolev 空间的自反性提供弱极限，闭凸约束保留边界数据，弱下半连续性完成能量比较。所获极小元首先属于弱函数空间；它是否更光滑，需要另外的方程估计。

\ProseParagraphBreak

山路方法比较的是路径上的最高能量。若这一高度附近没有临界点，沿负梯度方向作适当截断的下降流，就能把整条路径压低，违背高度的极小性。Palais--Smale 条件把能量受控且导数趋零的序列转成真正的临界点；次临界 Sobolev 紧嵌入则在半线性模型中提供所需的强收敛。下一章将以弱形式直接求解一般双线性问题，尤其处理不能来自对称能量的算子。
\end{chaptersummary}


\BookExercises

以下习题均为本书自拟。题目通过改变空间、边界或能量符号，检验变分结论的适用范围；最后两题计算山路高度及其随方程参数的变化。一维 Sobolev 函数的端点值按第一卷命题21.4.1取闭区间上的绝对连续代表，因此 \(u(1)-u(0)=\int_0^1u'\)。

% 自拟：无限维正定第二变分不足以保证极小。
\Needspace{6\baselineskip}
\begin{exercise}\label{b2:ex:14-positive-hessian}
在实空间 \(\ell^2\) 上令 \(J(x)=\sum_{n\ge1}x_n^2/n-\lVert x\rVert^4\)。证明 \(J\) 二阶连续 Fréchet 可微，零点为驻点且 \(D^2J(0)[h,h]>0\) 对每个 \(h\ne0\) 成立，但零点不是局部极小点。指出第二变分判别的哪项条件不成立。
\end{exercise}
\begin{solution}
有界自伴算子 \(Ax=(x_n/n)\) 给出 \(J(x)=\langle Ax,x\rangle-\langle x,x\rangle^2\)，所以它是连续多项式，特别是二阶连续可微。\(DJ(0)=0\)，且 \(D^2J(0)[h,h]=2\sum h_n^2/n>0\)。令 \(x^{(n)}=\sqrt{2/n}\,e_n\)，则 \(\lVert x^{(n)}\rVert\to0\)，而 \(J(x^{(n)})=2/n^2-4/n^2=-2/n^2<0=J(0)\)。在单位向量 \(e_n\) 上第二变分等于 \(2/n\to0\)，故不存在统一正下界 \(c\lVert h\rVert^2\)。
\end{solution}

% 自拟：带权端点问题，应用驻点和严格凸性。
\Needspace{6\baselineskip}
\begin{exercise}\label{b2:ex:14-weighted-string}
设 \(a_0>0\)、\(a\in C^1([0,1];\mathbb R)\)、\(a(x)\geq a_0\)。求 \(\tfrac12\int_0^1a(x)|u'(x)|^2\,\mathrm dx\) 在实函数 \(u\in H^1(0,1)\)、\(u(0)=0,u(1)=1\) 中的极小元和极小值，并用能量差证明唯一性。
\end{exercise}
\begin{solution}
记 \(L=\int_0^1a(s)^{-1}\,\mathrm ds\)，令 \(u(x)=L^{-1}\int_0^x a(s)^{-1}\,\mathrm ds\)，则 \(au'=L^{-1}\)。任一零端点扰动 \(h\in H_0^1(0,1)\) 满足 \(\int au'h'=L^{-1}(h(1)-h(0))=0\)，此等式也可由光滑逼近得到。因此能量差为 \(\tfrac12\int a|h'|^2\ge0\)，且等号迫使 \(h'=0\)、\(h=0\)。极小值是 \((2L^2)^{-1}\int a^{-1}=1/(2L)\)。
\end{solution}

% 自拟：边界弹簧与外法向符号。
\Needspace{6\baselineskip}
\begin{exercise}\label{b2:ex:14-boundary-spring}
给定 \(\beta>0\)、\(q\in\mathbb R\)，在实函数 \(u\in H^1(0,1)\)、\(u(0)=0\) 中极小化
\[J(u)=\tfrac12\int_0^1|u'|^2\,\mathrm dx+\tfrac\beta2u(1)^2-qu(1).\]
求极小元，并核验自然边界条件与全局极小性。
\end{exercise}
\begin{solution}
第一变分为 \(\int u'h'+(\beta u(1)-q)h(1)\)，其中 \(h(0)=0\)。内部方程是 \(u''=0\)，右端条件是 \(u'(1)+\beta u(1)=q\)，所以 \(u(x)=qx/(1+\beta)\)。对任意允许扰动，交叉项消失，能量差为 \(\tfrac12\int|h'|^2+\tfrac\beta2h(1)^2\)。它非负，等号时 \(h=0\)，故所得解是唯一全局极小元。
\end{solution}

% 自拟：积分约束产生常数乘子，含显式解。
\Needspace{6\baselineskip}
\begin{exercise}\label{b2:ex:14-mass-constraint}
给定 \(M\in\mathbb R\)，在实空间 \(H_0^1(0,1;\mathbb R)\) 中，附加 \(\int_0^1u=M\)，极小化 \(\tfrac12\int_0^1|u'|^2\)。不用抽象乘子定理，求出唯一极小元和极小值。
\end{exercise}
\begin{solution}
取 \(u(x)=6Mx(1-x)\)，则端点为零且积分为 \(M\)。允许扰动 \(h\) 满足零端点和零积分；分部积分给出 \(\int u'h'=-\int u''h=12M\int h=0\)。能量差为 \(\tfrac12\lVert h'\rVert_2^2\)，故它是唯一极小元。又 \(\int_0^1(1-2x)^2\,\mathrm dx=1/3\)，极小值为 \(6M^2\)。常数 \(-u''=12M\) 正是积分约束所允许留下的乘子；若错误地令所有紧支集方向都可行，就会误得 \(u''=0\)。
\end{solution}

% 自拟：临界负势能与最优Poincare常数。
\Needspace{6\baselineskip}
\begin{exercise}\label{b2:ex:14-negative-potential}
设 \(\lambda\in\mathbb R\)，在实空间 \(H_0^1(0,\pi;\mathbb R)\) 上考察 \(J_\lambda(u)=\tfrac12\int_0^\pi(|u'|^2-\lambda|u|^2)\)。使用区间正弦展开所得 \(\lVert u\rVert_2\le\lVert u'\rVert_2\)，分别讨论 \(\lambda<1\)、\(\lambda=1\)、\(\lambda>1\) 的极小元与强制性。
\end{exercise}
\begin{solution}
若 \(\lambda<1\)，当 \(0\le\lambda<1\) 时 \(J_\lambda\ge(1-\lambda)\lVert u'\rVert_2^2/2\)，当 \(\lambda<0\) 时更有 \(J_\lambda\ge\lVert u'\rVert_2^2/2\)，故强制且唯一极小元为零。若 \(\lambda=1\)，正弦展开给出 \(2J_1=\frac\pi2\sum_{n\ge1}(n^2-1)|a_n|^2\)，所以极小元恰为 \(c\sin x\)，泛函不强制。若 \(\lambda>1\)，沿 \(u=c\sin x\) 有 \(J_\lambda=\pi(1-\lambda)c^2/4\to-\infty\)，故没有有限极小值。展开公式先对有限正弦和成立，再由 \(H_0^1\) 逼近推广。
\end{solution}

% 自拟：非对称双线性型的有限维反例。
\Needspace{6\baselineskip}
\begin{exercise}\label{b2:ex:14-skew-part}
令 \(A=\left(\begin{smallmatrix}2&-3\\3&2\end{smallmatrix}\right)\)、\(f=(1,0)\)。对 \(u\in\mathbb R^2\)，分别求 \(Au=f\) 的解及 \(\tfrac12u^{\mathsf T}Au-f\cdot u\) 的极小元。证明 \((Au)\cdot u\ge2|u|^2\)，并解释这一强制性为何不足以使两个问题相同。
\end{exercise}
\begin{solution}
\(A^{-1}=13^{-1}\left(\begin{smallmatrix}2&3\\-3&2\end{smallmatrix}\right)\)，故方程解为 \((2/13,-3/13)\)。由于 \(u^{\mathsf T}Au=2|u|^2\)，泛函是 \(|u|^2-u_1\)，唯一极小元为 \((1/2,0)\)。强制性等式确实成立，但二次型只保留 \((A+A^{\mathsf T})/2=2I\)；驻点方程为 \(2u=f\)，反对称部分没有进入能量。
\end{solution}

% 自拟：障碍约束产生变分不等式。
\Needspace{6\baselineskip}
\begin{exercise}\label{b2:ex:14-obstacle}
设整数 \(d\geq1\)，\(\Omega\subset\mathbb R^d\) 为非空有界开集，\(K\) 是实空间 \(H_0^1(\Omega;\mathbb R)\) 的非空范数闭凸集，\(f\in L^2(\Omega;\mathbb R)\)。令 \(E(v)=\frac12\int_\Omega|\nabla v|^2-\int_\Omega fv\)。证明 \(E\) 在 \(K\) 上有唯一极小元 \(u\)，并证明它等价于
\[\int_\Omega\nabla u\cdot\nabla(v-u)\ge\int_\Omega f(v-u)\qquad(v\in K).\]
\end{exercise}
\begin{solution}
Poincaré 不等式给出强制性；闭凸集弱闭，能量凸且范数连续，故弱下半连续。第一卷定理 9.5.9 给出存在性，梯度平方的严格凸性给出唯一性。若 \(u\) 极小，对 \(u+t(v-u)\in K\)、\(0<t\le1\) 展开能量，除以 \(t\) 令 \(t\downarrow0\)，即得不等式。反之，取 \(h=v-u\) 于精确二次展开，线性项由不等式非负，平方项也非负，故 \(E(v)\ge E(u)\)。凸约束只提供单侧可行线段，因而得到不等式。
\end{solution}

% 自拟：受迫阻尼波的能量估计。
\Needspace{6\baselineskip}
\begin{exercise}\label{b2:ex:14-damped-wave}
设整数 \(d\geq1\)，\(c,T,\gamma>0\)，\(\Omega\subset\mathbb R^d\) 为非空有界 \(C^1\) 开集，实函数 \(u\) 在 \([0,T]\times\overline\Omega\) 的某开邻域内属于 \(C^2\)。设 \(f\in C([0,T]\times\overline\Omega;\mathbb R)\)，且 \(u_{tt}-c^2\Delta u+\gamma u_t=f\) 于 \((0,T)\times\Omega\)，\(u(t,y)=0\) 对 \(0\leq t\leq T\)、\(y\in\partial\Omega\) 成立。记
\[
 \mathcal E(t)=\frac12\int_\Omega
            (|u_t(t,x)|^2+c^2|\nabla u(t,x)|^2)\,\mathrm{d}x.
\]
证明对 \(0\le t\le T\)，
\[\mathcal E(t)+\frac\gamma2\int_0^t\lVert u_s\rVert_2^2\,\mathrm ds
 \le\mathcal E(0)+\frac1{2\gamma}\int_0^t\lVert f(s)\rVert_2^2\,\mathrm ds.\]
说明该估计怎样给出相同初边值数据的唯一性。
\end{exercise}
\begin{solution}
乘方程以 \(u_t\) 并分部积分，得 \(\mathcal E'+\gamma\lVert u_t\rVert_2^2=\int fu_t\)。由 Young 不等式，右端不超过 \(\lVert f\rVert_2^2/(2\gamma)+\gamma\lVert u_t\rVert_2^2/2\)。移项并对时间积分，得到所需估计。对两个解之差应用，右端外力与初始能量都为零，故差的时间导数和梯度均为零；零边界或零初位移固定剩余常数，差恒为零。
\end{solution}

% 自拟：一维p能量、端点约束与非线性通量；常外力部分增加实质训练，
% 正文存在性、可微性和严格凸性的论证均不依赖本题。
\Needspace{6\baselineskip}
\begin{exercise}\label{b2:ex:14-p-laplace-interval}
设 \(1<p<\infty\)，\(p'=p/(p-1)\)，
\[
 \mathcal C=\{u\in W^{1,p}(0,1;\mathbb R):u(0)=0,\ u(1)=1\},
 \qquad J_p(u)=\frac1p\int_0^1|u'|^p\,\mathrm{d}x.
\]
端点值取绝对连续代表。令 \(A_p(s)=|s|^{p-2}s\)（\(s\ne0\)），并规定 \(A_p(0)=0\)。
\begin{enumerate}
\item\label{b2:item:14-p-interval-unforced}证明 \(J_p\) 在 \(\mathcal C\) 上的唯一极小元为 \(u(x)=x\)，并求极小值。由第一变分写出
\(\int_0^1 A_p(u')h'=0\)（\(h\in W_0^{1,p}(0,1;\mathbb R)\)），
解释这就是分布方程 \((A_p(u'))'=0\)，并指出严格凸性如何保证唯一性。
\item\label{b2:item:14-p-interval-forced}给定 \(\lambda\in\mathbb R\)，在同一类 \(\mathcal C\) 上极小化
\[
 J_{p,\lambda}(u)=J_p(u)-\lambda\int_0^1u(x)\,\mathrm{d}x.
\]
写出以一个实常数 \(C\) 为参数的候选解，并证明端点条件唯一确定 \(C\)，所得候选确为唯一全局极小元。所用方程应写成非线性通量的弱形式，不预设 \(u\) 具有经典二阶导数。
\end{enumerate}
\end{exercise}
\begin{solution}
对 \(u\in\mathcal C\)，绝对连续表示给出 \(\int_0^1u'=1\)。
函数 \(s\mapsto|s|^p\) 严格凸，Jensen 不等式因而给出
\[
 J_p(u)\geq\frac1p\left|\int_0^1u'\right|^p=\frac1p.
\]
等号成立当且仅当 \(u'\) 几乎处处为常数；该常数只能为 \(1\)，再由 \(u(0)=0\)
得到 \(u(x)=x\)。这同时给出极小值与唯一性，而非只验证一个驻点。

允许方向恰为 \(h\in W_0^{1,p}(0,1;\mathbb R)\)。
由\cref{b2:prop:14-power-frechet}，极小点满足
\[
 DJ_p(u)h=\int_0^1A_p(u')h'\,\mathrm{d}x=0.
\]
特别取 \(h\in C_c^\infty(0,1)\)，便是所列分布方程。
反过来，若通量 \(A_p(u')\in L^{p'}(0,1)\) 的分布导数为零，由一维绝对连续代表的积分表示，它几乎处处为常数；
\(A_p\) 严格递增且可逆，故 \(u'\) 也为常数，端点再将它固定为 \(1\)。
这里凸性把驻点变成全局极小点，严格凸性再排除不同的极小点。

对常外力问题，第一变分为
\[
 \int_0^1 A_p(u')h'\,\mathrm{d}x
       =\lambda\int_0^1h\,\mathrm{d}x
 \qquad(h\in W_0^{1,p}(0,1;\mathbb R)),
\]
即 \(-(A_p(u'))'=\lambda\)。
记 \(B_p(q)=\operatorname{sgn}(q)|q|^{1/(p-1)}\)，其中 \(B_p(0)=0\)；
这是 \(A_p\) 的连续严格递增逆映射。积分通量方程得到候选
\[
 u_C(x)=\int_0^x B_p(C-\lambda s)\,\mathrm{d}s.
\]
它属于 \(W^{1,p}(0,1)\)，满足左端值；右端值要求
\[
 F_\lambda(C)\coloneqq\int_0^1 B_p(C-\lambda s)\,\mathrm{d}s=1.
\]
在任何紧的 \(C\) 区间内，被积函数统一有界且连续，故 \(F_\lambda\) 连续。
\(C\) 严格增大时每个被积值都严格增大，所以 \(F_\lambda\) 严格递增。
又 \(C-|\lambda|\leq C-\lambda s\leq C+|\lambda|\)，
由 \(B_p\) 的单调性可知 \(F_\lambda(C)\to\pm\infty\) 当 \(C\to\pm\infty\)。
介值定理于是给出唯一的 \(C\)。

所构造的 \(u=u_C\) 满足 \(A_p(u')=C-\lambda x\)。
对绝对连续且零端点的 \(h\)，分部积分直接给出上述弱等式；
这里微分的是通量，不需要先求 \(u''\)。
最后取任意 \(v\in\mathcal C\)，由凸函数的切线不等式及弱等式，
\[
 \begin{aligned}
 J_{p,\lambda}(v)-J_{p,\lambda}(u)
 &\geq\int_0^1A_p(u')(v'-u')\,\mathrm{d}x
          -\lambda\int_0^1(v-u)\,\mathrm{d}x\\
 &=0.
 \end{aligned}
\]
所以候选为全局极小元。严格凸性使等号必须满足 \(v'=u'\) 几乎处处，
再由同一端点值得 \(v=u\)。当 \(\lambda=0\) 时，参数为 \(C=1\)，公式退回 \(u(x)=x\)。
\end{solution}

% 自拟选读题：二维临界点与路径障碍；增加临界点类型和极小极大值的直接计算。
\Needspace{6\baselineskip}
\begin{optionalexercise}[一个可以直接计算的山口]\label{b2:ex:14-mountain-finite}
在实空间 \(\mathbb R^2\) 上令
\[
 J(x,y)=\frac{x^2+y^2}{2}-\frac{x^4}{4},
 \quad e=(2,0),
 \quad\Gamma=\{\gamma\in C([0,1];\mathbb R^2):\gamma(0)=(0,0),\ \gamma(1)=e\}.
\]
不用抽象山路定理，计算 \(c=\inf_{\gamma\in\Gamma}\max_tJ(\gamma(t))\)。求出全部临界点，利用二阶导数区分它们的类型，并证明 \(J\) 满足 Palais–Smale 条件。
\end{optionalexercise}
\begin{solution}
每条路径的第一坐标都从 \(0\) 走到 \(2\)，因而在某点等于 \(1\)。该点的能量为 \(1/4+y^2/2\ge1/4\)，故 \(c\ge1/4\)。线段路径 \(\gamma(t)=(2t,0)\) 的最大能量等于 \(1/4\)，所以 \(c=1/4\)。

梯度为 \((x-x^3,y)\)，临界点恰为 \((0,0)\)、\((1,0)\)、\((-1,0)\)。Hessian 为 \(\operatorname{diag}(1-3x^2,1)\)。它在原点正定，给出严格局部极小；在另外两点有一个负特征值和一个正特征值，给出鞍点。若 \(\nabla J(x_n,y_n)\to0\)，则 \(y_n\to0\) 且 \(x_n-x_n^3\to0\)。后者迫使 \((x_n)\) 有界，否则三次项的绝对值趋于无穷。因此序列在有限维空间中有收敛子列，特别满足 Palais–Smale 条件。对这个多项式，梯度趋零本身就迫使序列有界。
\end{solution}

% 自拟选读题：参数缩放和路径族双射；不复述正文PS验证。
\Needspace{6\baselineskip}
\begin{optionalexercise}[非线性强度与山路高度]\label{b2:ex:14-mountain-scaling}
设整数 \(d\ge1\)，\(\Omega\subset\mathbb R^d\) 为非空有界开集，\(2<p<2^*\)，其中 \(2^*=2d/(d-2)\) 当 \(d\ge3\)，\(2^*=+\infty\) 当 \(d=1,2\)。在实空间 \(H=H_0^1(\Omega)\) 上令
\[
 J_\lambda(u)=\frac12\|\nabla u\|_2^2-\frac\lambda p\|u\|_p^p
 \quad(\lambda>0).
\]
固定 \(e_1\in H\) 使 \(J_1(e_1)<0\)，令 \(e_\lambda=\lambda^{-1/(p-2)}e_1\)，并以从 \(0\) 到 \(e_\lambda\) 的所有连续路径定义山路高度 \(c_\lambda\)。证明
\[
 c_\lambda=\lambda^{-2/(p-2)}c_1.
\]
若 \(u_1\ne0\) 是 \(-\Delta u=|u|^{p-2}u\) 的零边界弱解，验证 \(u_\lambda=\lambda^{-1/(p-2)}u_1\) 是 \(-\Delta u=\lambda|u|^{p-2}u\) 的非零弱解。解释为什么比较山路高度时，端点也需要随参数缩放。
\end{optionalexercise}
\begin{solution}
设 \(s=\lambda^{-1/(p-2)}>0\)。则 \(\lambda s^p=s^2\)，所以对每个 \(v\in H\) 有 \(J_\lambda(sv)=s^2J_1(v)\)。映射 \(\gamma\mapsto s\gamma\) 将从 \(0\) 到 \(e_1\) 的路径族双射到从 \(0\) 到 \(e_\lambda\) 的路径族，其逆为除以 \(s\)。在每条路径上先取最大值，再对双射的两个路径族取下确界，即得 \(c_\lambda=s^2c_1\)。

对任意 \(v\in H\)，原弱方程给出
\[
 \int\nabla(su_1)\cdot\nabla v
 =s\int|u_1|^{p-2}u_1v
 =\lambda s^{p-1}\int|u_1|^{p-2}u_1v,
\]
最后一步使用 \(\lambda s^{p-2}=1\)，这正是所求弱方程。因 \(s>0\)，非零性保留。若强行保持同一端点，缩放后的路径通常不属于待比较的路径族，上面的双射就不成立；因此能量缩放公式本身不足以推出那个不同路径族的高度公式。
\end{solution}
